\chapter{Composite kernels and the odd two-body kernels}

Chapter 10 wrote the dilatation component $D_1$ and the translation component $P_1$ of a local operator as four times the logarithmic pair kernels, a one-body term and a composite term, the graphs in which $S^{(1)}(\Delta)$ or $S^{(1)}(P_b)$ sits on the cycle as a line from a ghost leg to a gauge leg (Theorem \ref{thm:c10-D1}, Proposition \ref{prop:c10-P1}), and evaluated the composite term on the planar scalar sector. On a general operator the locality of the composite term is a statement to be proved: on the letters $A|_V$, $\lambda$, $c$ the renormalised composite two-point functional contains $\log|p|^2$ of the momentum (Remark \ref{rem:c9-Pi}). This chapter shows that $Q$ maps the composite graphs of a cocycle to cycles with one special edge, so that the composite term is a sum over the pairs of letters of a local kernel, the composite pair kernel; it computes these kernels as simplex integrals and evaluates them on twist two, on the $\mathfrak{su}(1|1)$ sector, on the fermion and gluon $\mathfrak{sl}(2)$ sectors and on two mixed families. On all of them the composite pair kernels satisfy the two identities that Corollary \ref{cor:c10-P7} requires, and with the arity-two equation along $(\Delta,P_b)$ they fix the three sector coefficients of $\hbar F_1^{(1)}(\Delta)$ to $-\tfrac13$ units: $D_1$ is four times the map of Theorem \ref{thm:c10-brst}, with spectrum twice the one-loop anomalous dimensions, and $P_1$ vanishes on cohomology in the translation-covariant scheme. The last section sets up the two-body kernels of $Q_1$, $S_1$ and $K_1$ in the same scheme: their graphs and output bounds, the relations of $S_1$ to $Q_1$ and of $K_1$ to $P_1$, their renormalisation freedom and the comparison with the one-loop differential of the holomorphic twist.

The unit is $\frac{\hbar\operatorname{tr}_{\mathfrak g}(\mathrm{ad}_X\mathrm{ad}_Y)}{16\pi^2(X,Y)}$, equal to $-\frac{g^2N}{32\pi^2}$ for $\mathfrak g=\mathfrak{su}(N)$ (Remark \ref{rem:c6-konishi}). Edges carry $P(\varepsilon,T)$ of Definition \ref{def:c8-PT}; the special edge is $K_\varepsilon-K_0=-\int_0^\varepsilon(H\otimes1)K_t\,dt$, the factor $e^{-\varepsilon|p|^2}-1$ times the components of $K_0$ (Proposition \ref{prop:c4-defect1}). Every component below is independent of the endpoint $T\in(0,\infty]$ of the transfer (Theorem \ref{thm:c8-L} (ii)) and is evaluated at $T=\infty$, where the graph sums converge at generic momenta (Proposition \ref{prop:c8-Tsafe}); along $K_b$ it is taken at a finite $T$ (Definition \ref{def:c8-K}). On arity zero $s=Q+\{I,-\}$ (Theorem \ref{thm:c8-obstruction}). Letters, the pair weight, the pair degree and the Koszul signs are those of Definition \ref{def:c10-letters} and \eqref{eq:c10-decomp}. The field-counting functional is $2\int(A,A^*)-4\int(\lambda^*,\lambda)-\int(c^*,c)$, whose bracket multiplies the coordinate functionals of $c$, $A$, $\lambda$, $A^*$, $\lambda^*$, $c^*$ by $1$, $2$, $4$, $-2$, $-4$, $-1$ (Lemma \ref{lem:c6-count}), and the map of Theorem \ref{thm:c10-brst} is the coefficient of $\log\varepsilon$ in $\partial_\delta I^{CT}[\varepsilon](O)$ minus the unit times the bracket of $O$ with it (Section 10.7). Along $\Delta$, $v=2E$ and $v.$ is the operator of Lemma \ref{lem:c9-SDelta}, so that $\{S^{(1)}(\Delta),M\}=v.M$; along $P_b$, $v.$ is $\partial_b$ on every slot (Proposition \ref{prop:c7-cov}); for $x\in\{\Delta,P_b\}$, $v_x.$ is the operator of the vector field of $x$ (Definition \ref{def:c7-vdot}). The two-component form of Definition \ref{def:c8-twocomp} of every main formula is in Remark \ref{rem:c11-twocomp}.

\section{The composite term as a sum over pairs}

\begin{definition}[Square kernel and composite kernels]
\label{def:c11-kernels}
For $0<\varepsilon<T\le\infty$ and $P=P(\varepsilon,T)$, the square kernel is $\big((P\star)^2\otimes1\big)K_0$, and the composite kernel along $x\in\{\Delta,P_b\}$ is
\[
\big(P\star\,v_x.\otimes1+1\otimes P\star\,v_x.\big)P,
\]
the two legs of $S^{(1)}(x)$ joined, in both orders, to two edges $P$ (Definition \ref{def:c9-composite}, Theorem \ref{thm:c9-P1}). The sector parts of $S^{(1)}(P_b)$ are
\[
\int\big(\partial_b(A|_{W_{\mathbb R}^\perp}),A^*|_{W_{\mathbb R}^\perp}\big),\qquad -\int(\partial_b\lambda,\lambda^*),\qquad \int\big(\partial_b(A|_V),A^*|_V\big)-\int(\partial_bc,c^*);
\]
their sum is $S^{(1)}(P_b)$, and each sector part of $S^{(1)}(\Delta)$ of Section 10.8 is $2\int x^\mu$ times the density of the sector part of $S^{(1)}(P_{e_\mu})$ of the same sector, summed over $\mu$, plus a translation-invariant functional, with $e_\mu$ an orthonormal basis of $V$ and $x^\mu$ the coordinates.
\end{definition}

\begin{lemma}[Structure of the composite kernels]
\label{lem:c11-kernels}
\begin{itemize}
\item (i) The composite kernel along $x$ is $2\big(P\star\,v_x.\,P\star\otimes1\big)K_0=(v_x.\otimes1-1\otimes v_x.)$ applied to the square kernel.
\item (ii) $(v_x.\otimes1+1\otimes v_x.)$ applied to the square kernel is $-4\big(\varepsilon\partial_\varepsilon+T\partial_T\big)$ of the square kernel for $x=\Delta$ and $0$ for $x=P_b$. Hence the composite kernel along $\Delta$ is $2(v.\otimes1)$ of the square kernel plus $4(\varepsilon\partial_\varepsilon+T\partial_T)$ of it, and the composite kernel along $P_b$ is $2(\partial_b\otimes1)$ of the square kernel.
\item (iii) The square kernel and the composite kernels have components in $c\otimes A|_V$ and $A|_V\otimes c$ alone: as operators, $(P\star)^2$ and $P\star\,v_x.\,P\star$ map the slot of $A^*|_V$ to the slot of $c$ and the slot of $c^*$ to the slot of $A|_V$, and vanish on the other slots.
\item (iv) In a graph the composite kernel is the vertex $S^{(1)}(x)$ with both legs on edges $P$. Its term $\int(v_x.A|_V,A^*|_V)$ contributes, with the leg $A|_V$ on the component $A\otimes A$ of $P$ and the leg $A^*|_V$ on the component $A^*|_V\otimes c$. Along $\Delta$ this internal vertex is
\[
\int\big((2E+2)A|_V,A^*|_V\big)=2\int x^\mu\big(\partial_\mu(A|_V),A^*|_V\big)+2\int\big(A|_V,A^*|_V\big),
\]
along $P_b$ it is $\int(\partial_b(A|_V),A^*|_V)$. The components of $P$ lie in $A\otimes A$, $\lambda\otimes\lambda$, $c\otimes A^*|_V$ and $A^*|_V\otimes c$ (Proposition \ref{prop:c3-P}), so the terms of $S^{(1)}(x)$ in the sectors of $A|_{W_{\mathbb R}^\perp}$, $\lambda$ and $c$, whose antifield legs are $A^*|_{W_{\mathbb R}^\perp}$, $\lambda^*$, $c^*$, enter a cycle through the special edge.
\item (v) $(Q\otimes1+1\otimes Q)$ annihilates the composite kernel; the operator $\partial$ of the composite kernel is odd and anticommutes with $Q$.
\item (vi) At $T=\infty$ the component of the square kernel is $(Q^{GF})^2$ on the slot of $A^*|_V$, which is $2d^*$ with values in the slot of $c$ (Proposition \ref{prop:c9-composite} (ii)), times $(-,-)^{-1}$ and
\[
\int_\varepsilon^\infty\!\!\int_\varepsilon^\infty e^{-(s'+s'')|p|^2}\,ds'\,ds''=e^{-2\varepsilon|p|^2}|p|^{-4}:
\]
a doubled edge with two Schwinger times $s',s''\in[\varepsilon,\infty)$, of momentum degree $-3$, which obeys the weight rule of Proposition \ref{prop:c3-P} for the slots $c$ (weight $0$) and $A|_V$ (weight $1$).
\end{itemize}
\end{lemma}

\begin{proof}
(i) The first equality is Proposition \ref{prop:c9-composite} (i) along $\Delta$; along $P_b$ the same computation holds, $\partial_b$ preserving the pairing and commuting with $Q^{GF}$ and $H$ (Proposition \ref{prop:c7-cov}). For the second, $2P\star\,v_x.\,P\star=v_x.(P\star)^2+(P\star)^2v_x.$, along $\Delta$ because $[v.,P\star]$ is $Q^{GF}$ times a function of $H$ and commutes with $P\star$ (proof of Proposition \ref{prop:c9-composite} (i)), along $P_b$ because $\partial_b$ commutes with $P\star$. For the even operator $X=(P\star)^2$, $(X\,v_x.\otimes1)K_0=(X\otimes1)(v_x.\otimes1)K_0=-(1\otimes v_x.)(X\otimes1)K_0$ by Lemma \ref{lem:c9-ops}, $v_x.$ being even and pairing preserving.

(ii) Since $(v_x.\otimes1+1\otimes v_x.)K_0=0$, the left side is $([v_x.,X]\otimes1)K_0$. Along $\Delta$, $X=(Q^{GF})^2g(H)^2$ with $g(H)=\int_\varepsilon^Te^{-tH}dt$, $[v.,Q^{GF}]=-4Q^{GF}$ and $[v.,f(H)]=-4Hf'(H)$ (Proposition \ref{prop:c7-dil}). From $Hg'(H)=Te^{-TH}-\varepsilon e^{-\varepsilon H}-g=(\varepsilon\partial_\varepsilon+T\partial_T)g-g$,
\[
[v.,X]=-8X-4(Q^{GF})^2H\partial_H\big(g^2\big)=-8X-4\big((\varepsilon\partial_\varepsilon+T\partial_T)X-2X\big)=-4\big(\varepsilon\partial_\varepsilon+T\partial_T\big)X .
\]
Along $P_b$, $[\partial_b,X]=0$. Insert into (i).

(iii) $(Q^{GF})^2$ vanishes on the slots of $c$, $A$, $\lambda$, $\lambda^*$; on $A^*$ and $c^*$ it is $2d^*(A^*|_V)$ in the slot of $c$ and $-2dc^*$ in the slot of $A|_V$ (Proposition \ref{prop:c9-composite} (ii)); $v_x.$ preserves the slots, and the components $A^*\otimes A$ and $c^*\otimes c$ of $K_0$ go to $c\otimes A|_V$ and $A|_V\otimes c$.

(iv) Applied to $A^*\otimes A$, the kernel $(P\star\,v_x.\,P\star\otimes1)K_0$ is the component $A^*\to A$ of $Q^{GF}$, then $v_x.$ on the slot of $A$, then the component $A\to c$ of $Q^{GF}$: the vertex $\int(v_x.A|_V,A^*|_V)$ between an $A\otimes A$ and an $A^*|_V\otimes c$ component, the factor $2$ counting the two orientations. On $A|_V$, $L_{2E}A|_V=2x^\mu\partial_\mu(A|_V)+2A|_V$ (Definition \ref{def:c7-vdot}).

(v) Proposition \ref{prop:c9-composite} (iii) along $\Delta$; along $P_b$, $QP\star+P\star Q=e^{-\varepsilon H}-e^{-TH}$ commutes with $\partial_b$, which commutes with $Q$, so $Q(P\star\,\partial_b\,P\star)=(P\star\,\partial_b\,P\star)Q$. The identity $[\partial_\phi,Q]=\partial_{(Q\otimes1+1\otimes Q)\phi}$ of Lemma \ref{lem:c9-ops} gives the anticommutation; the composite kernel has degree $-1$.

(vi) Proposition \ref{prop:c9-composite} (iv) and $\big(\int_\varepsilon^\infty e^{-s|p|^2}ds\big)^2=e^{-2\varepsilon|p|^2}|p|^{-4}$.
\end{proof}

\begin{definition}[Composite graphs and the arity-one counterterm of an operator]
\label{def:c11-graphs}
Let $M$ be a local functional, $\delta$ a square-zero parameter of degree $-|M|$, $x\in\{\Delta,P_b\}$ and $P=P(\varepsilon,\infty)$. The composite graphs of $M$ along $x$ are the $\delta$-linear part of
\[
\tfrac\hbar2\,\partial_{\text{composite kernel along }x}\,W(P,I+\delta M)_{\rm tree}:
\]
the composite kernel closes two legs of a tree of vertices $I$ and one vertex $M$. The cycle so formed either passes through two letters $X_a$, $X_b$ of $M$, a composite cycle through the pair $(X_a,X_b)$ with trees of $I$ on its outputs, or passes through vertices of $I$ alone and is joined to $M$ by a tree through one letter, a composite loop of $I$ on a leg of $M$. The arity-one counterterm of $M$ along $x$ is $\hbar\,\mathrm{Sing}_\varepsilon$ of the amputated composite cycles through the pairs, summed over the pairs with the Koszul signs of \eqref{eq:c10-decomp}. The singular part of the composite loops of $I$ is $I^{CT}_{\mathfrak{psl}(4|4)}[\varepsilon]^{(1)}(x)$ (Proposition \ref{prop:c9-structure} (i)); its quadratic part along $\Delta$ is the unit times $\log\varepsilon$ times $Q\big(\int(\iota_xA|_V,c^*)-\tfrac12\int(\rho(x)\lambda^*,\lambda^*)\big)$, with no term in $\varepsilon^{-1}$ (Theorem \ref{thm:c9-quadratic} (i)).
\end{definition}

\begin{proposition}[The image of the composite graphs under $Q$]
\label{prop:c11-master}
For every local functional $M$, $x\in\{\Delta,P_b\}$ and $\varepsilon>0$, at $T=\infty$,
\[
Q\big(\text{composite graphs of }M\text{ along }x\big)=-\tfrac\hbar2\,\partial_{\text{composite kernel along }x}\Big[\text{tree-level flow, }\delta\text{-linear, with one vertex }sM\text{ or one special edge}\Big],
\]
the special edge joining a vertex $I$ to a letter of $M$, from $\{I,M\}_\varepsilon-\{I,M\}$, or two vertices $I$, from $\tfrac12(\{I,I\}_\varepsilon-\{I,I\})$. Since $\Delta_\varepsilon-\Delta_0=-\partial_{K_\varepsilon-K_0}$ (Definition \ref{def:c1-bracket}, Proposition \ref{prop:c4-defect1}), $\{f,g\}_\varepsilon-\{f,g\}$ is minus $f$ and $g$ joined by one edge $K_\varepsilon-K_0$; for $sM=0$ the right side is $\tfrac\hbar2\,\partial_{\text{composite kernel along }x}$ of the trees with exactly one edge $K_\varepsilon-K_0$.
\end{proposition}

\begin{proof}
By Lemma \ref{lem:c11-kernels} (v), $Q$ anticommutes with $\partial$ of the composite kernel. By Proposition \ref{prop:c4-defectflow} at order $\hbar^0$, $Q\,W(P,J)_{\rm tree}$ is the tree-level flow of $QJ+\tfrac12\{J,J\}_\varepsilon$; the term $\hbar\Delta_\varepsilon J$ is of order $\hbar$, and its composite closure has two cycles. For $J=I+\delta M$ the classical master equation $QI+\tfrac12\{I,I\}=0$ (Proposition \ref{prop:c1-cme}) leaves $\tfrac12(\{I,I\}_\varepsilon-\{I,I\})$ in the $\delta$-independent part and $sM+(\{I,M\}_\varepsilon-\{I,M\})$ in the $\delta$-linear part.
\end{proof}

\begin{definition}[Special-edge composite cycles and the composite pair kernel]
\label{def:c11-sp}
For two letters $X_a$, $X_b$ of any slots and $x\in\{\Delta,P_b\}$, the special-edge composite cycles of $(X_a,X_b)$ along $x$ are
\[
\tfrac\hbar2\,\partial_{\text{composite kernel along }x}\,\partial_{K_\varepsilon-K_0}\big(\text{amputated chains from }X_a\text{ to }X_b\big):
\]
the chains of Definition \ref{def:c10-chain} with $k\ge1$ vertices of $I$ in which one edge $e_c$ is the composite kernel along $x$, one other edge $e_\sigma$ is the special edge, with the components of $K_0$ (each field slot with its antifield slot), which may end at a letter of any slot, and every remaining edge is $P(\varepsilon,\infty)$. By Proposition \ref{prop:c11-master} they are minus $\tfrac\hbar2\,\partial$ of the composite kernel applied to the chains with one bracket-defect vertex. A chain is read with the signs of Lemma \ref{lem:c10-sign} along the chain from $X_a$, times
\[
(-1)^{n(<\sigma)+n(<c)+[c<\sigma]},
\]
where $n(<j)$ is the number of odd outputs at the vertices before the edge $e_j$, and $[c<\sigma]$ is $1$ if $e_c$ precedes $e_\sigma$ and $0$ otherwise; the outputs are left open as in Definition \ref{def:c10-pair}. The composite pair kernel of $(X_a,X_b)$ along $x$ is the coefficient of $\varepsilon^0$ of its special-edge composite cycles.
\end{definition}

By Proposition \ref{prop:c11-master}, for $sM=0$ the image under $Q$ of the composite cycles through the pairs and of the composite loops on the legs is the sum of the special-edge composite cycles with trees of $I$ on their outputs, together with the composite closures in which the special edge lies on a tree; Theorem \ref{thm:c11-local} identifies the composite term of a cocycle with the sum over its pairs of letters of the composite pair kernels.

\begin{proposition}[Reachability and census]
\label{prop:c11-census}
\begin{itemize}
\item (i) The composite kernel joins a slot $c$ to a slot $A|_V$. The letters $c$ and $A|_V$ reach it directly; the letters $A|_{W_{\mathbb R}^\perp}$, $\lambda$ and $A^*|_V$ reach it through one vertex, $A|_V(A|_{W_{\mathbb R}^\perp})^2$ entered at $A|_{W_{\mathbb R}^\perp}$, $\lambda|_{S_+\otimes W}\lambda|_{S_-\otimes W^*}A|_V$ entered at $\lambda$, $([A|_V,c],A^*|_V)$ entered at $c$, each left at $A|_V$; the letters $A^*|_{W_{\mathbb R}^\perp}$, $\lambda^*$, $c^*$, on which $P$ and the composite kernel have no component, reach it through the special edge, to a slot $A|_{W_{\mathbb R}^\perp}$, $\lambda$, $c$.
\item (ii) A chain of Definition \ref{def:c11-sp} exists for $45$ of the $55$ unordered pairs of the ten slots; the other ten are the pairs of two letters among $A^*|_{W_{\mathbb R}^\perp}$, $\lambda^*|_{S_+\otimes W^*}$, $\lambda^*|_{S_-\otimes W}$, $c^*$, which would need two special edges. Every chain conserves the degree, the fermion parity and the centre $\mathbb Z/4$ of $SL(4,\mathbb C)$, and its outputs carry the pair degree; the composite kernel has degree $-1$ and the special edge $+1$.
\item (iii) Every chain has $k\ge1$ vertices. A chain with $k=1$ has the composite kernel at one letter and the special edge at the other, so $k=1$ occurs exactly for the pairs with a letter $A|_V$ and for the pairs $(c,A^*|_V)$, $(c,A^*|_{W_{\mathbb R}^\perp})$, $(c,\lambda^*)$, $(c,c^*)$; the other pairs, the field pairs $(A|_{W_{\mathbb R}^\perp},A|_{W_{\mathbb R}^\perp})$, $(A|_{W_{\mathbb R}^\perp},\lambda)$, $(\lambda,\lambda)$ among them, start at $k=2$.
\item (iv) At the pair weight without derivatives the leg sets of the ten pairs of field letters are those of Table \ref{tab:c11-fieldpairs}; the numbers of leg sets with $n\le3$ derivatives on the two letters, for the $21$ pairs of the active slots of Lemma \ref{lem:c10-active}, are those of Table \ref{tab:c11-counts}.
\end{itemize}
\end{proposition}

\begin{proof}
(i) Lemma \ref{lem:c11-kernels} (iii), Lemma \ref{lem:c10-active} and the vertices of Definition \ref{def:c1-vertex}. (ii), (iv) Enumeration of the chains with one composite kernel and one special edge, the vertices of $I$ and the components of $P$ and $K_0$, up to six vertices and output weight four above the pair weight; every chain is checked for degree, parity and the centre. A pair of two letters among $A^*|_{W_{\mathbb R}^\perp}$, $\lambda^*$, $c^*$ needs an edge with a component on each letter; $K_0$ alone has one, and a chain has one special edge. (iii) A chain with $k=0$ would join the two letters by the composite kernel and the special edge at once; the rest is read off from the enumeration.
\end{proof}

\begin{tabular}{l|c|l}
pair & pair weight & leg sets of the special-edge composite cycles, fewest vertices in brackets; ${}^\dagger$: position part along $\Delta$ and part along $P_b$\\
\hline
$(A|_V,A|_V)$ & $2$ & $A|_V$ $(1)$, $A|_V^2$ $(2)$, $(A|_{W_{\mathbb R}^\perp})^2$ $(2)$, $A|_V^3{}^\dagger$ $(2)$, $A|_V(A|_{W_{\mathbb R}^\perp})^2{}^\dagger$ $(2)$, $A^*|_V\,c^\dagger$ $(2)$, $\lambda|_{S_-\otimes W^*}\lambda|_{S_+\otimes W}{}^\dagger$ $(2)$\\
$(A|_V,A|_{W_{\mathbb R}^\perp})$ & $2$ & $A|_{W_{\mathbb R}^\perp}$ $(1)$, $A|_V\,A|_{W_{\mathbb R}^\perp}$ $(2)$, $A|_V^2A|_{W_{\mathbb R}^\perp}{}^\dagger$ $(2)$, $(A|_{W_{\mathbb R}^\perp})^3{}^\dagger$ $(2)$, $A^*|_{W_{\mathbb R}^\perp}c^\dagger$ $(2)$, $(\lambda|_{S_-\otimes W^*})^2{}^\dagger$ $(2)$, $(\lambda|_{S_+\otimes W})^2{}^\dagger$ $(2)$\\
$(A|_V,\lambda|_{S_+\otimes W})$ & $\tfrac52$ & $\lambda|_{S_+\otimes W}$ $(1)$, $A|_V\,\lambda|_{S_+\otimes W}$ $(2)$, $A|_{W_{\mathbb R}^\perp}\lambda|_{S_-\otimes W^*}$ $(2)$, $c\,\lambda^*|_{S_-\otimes W}$ $(2)$, $A|_V^2\lambda|_{S_+\otimes W}{}^\dagger$ $(2)$, $A|_V\,A|_{W_{\mathbb R}^\perp}\lambda|_{S_-\otimes W^*}{}^\dagger$ $(3)$, $A|_V\,c\,\lambda^*|_{S_-\otimes W}{}^\dagger$ $(3)$, $(A|_{W_{\mathbb R}^\perp})^2\lambda|_{S_+\otimes W}{}^\dagger$ $(2)$, $A|_{W_{\mathbb R}^\perp}c\,\lambda^*|_{S_+\otimes W^*}{}^\dagger$ $(3)$\\
$(A|_V,\lambda|_{S_-\otimes W^*})$ & $\tfrac52$ & $\lambda|_{S_-\otimes W^*}$ $(1)$, $A|_V\,\lambda|_{S_-\otimes W^*}$ $(2)$, $A|_{W_{\mathbb R}^\perp}\lambda|_{S_+\otimes W}$ $(2)$, $c\,\lambda^*|_{S_+\otimes W^*}$ $(2)$, $A|_V^2\lambda|_{S_-\otimes W^*}{}^\dagger$ $(2)$, $A|_V\,A|_{W_{\mathbb R}^\perp}\lambda|_{S_+\otimes W}{}^\dagger$ $(3)$, $A|_V\,c\,\lambda^*|_{S_+\otimes W^*}{}^\dagger$ $(3)$, $(A|_{W_{\mathbb R}^\perp})^2\lambda|_{S_-\otimes W^*}{}^\dagger$ $(2)$, $A|_{W_{\mathbb R}^\perp}c\,\lambda^*|_{S_-\otimes W}{}^\dagger$ $(3)$\\
$(A|_{W_{\mathbb R}^\perp},A|_{W_{\mathbb R}^\perp})$ & $2$ & $(A|_{W_{\mathbb R}^\perp})^2$ $(2)$, $A|_V(A|_{W_{\mathbb R}^\perp})^2{}^\dagger$ $(2)$\\
$(A|_{W_{\mathbb R}^\perp},\lambda|_{S_+\otimes W})$ & $\tfrac52$ & $A|_{W_{\mathbb R}^\perp}\lambda|_{S_+\otimes W}$ $(2)$, $A|_V\,A|_{W_{\mathbb R}^\perp}\lambda|_{S_+\otimes W}{}^\dagger$ $(2)$, $(A|_{W_{\mathbb R}^\perp})^2\lambda|_{S_-\otimes W^*}{}^\dagger$ $(3)$, $A|_{W_{\mathbb R}^\perp}c\,\lambda^*|_{S_-\otimes W}{}^\dagger$ $(3)$\\
$(A|_{W_{\mathbb R}^\perp},\lambda|_{S_-\otimes W^*})$ & $\tfrac52$ & $A|_{W_{\mathbb R}^\perp}\lambda|_{S_-\otimes W^*}$ $(2)$, $A|_V\,A|_{W_{\mathbb R}^\perp}\lambda|_{S_-\otimes W^*}{}^\dagger$ $(2)$, $(A|_{W_{\mathbb R}^\perp})^2\lambda|_{S_+\otimes W}{}^\dagger$ $(3)$, $A|_{W_{\mathbb R}^\perp}c\,\lambda^*|_{S_+\otimes W^*}{}^\dagger$ $(3)$\\
$(\lambda|_{S_+\otimes W},\lambda|_{S_+\otimes W})$ & $3$ & $(\lambda|_{S_+\otimes W})^2$ $(2)$, $A|_V(\lambda|_{S_+\otimes W})^2{}^\dagger$ $(3)$, $A|_{W_{\mathbb R}^\perp}\lambda|_{S_-\otimes W^*}\lambda|_{S_+\otimes W}{}^\dagger$ $(3)$, $c\,\lambda|_{S_+\otimes W}\lambda^*|_{S_-\otimes W}{}^\dagger$ $(3)$\\
$(\lambda|_{S_+\otimes W},\lambda|_{S_-\otimes W^*})$ & $3$ & $\lambda|_{S_-\otimes W^*}\lambda|_{S_+\otimes W}$ $(2)$, $A|_V\,\lambda|_{S_-\otimes W^*}\lambda|_{S_+\otimes W}{}^\dagger$ $(3)$, $A|_{W_{\mathbb R}^\perp}(\lambda|_{S_-\otimes W^*})^2{}^\dagger$ $(3)$, $A|_{W_{\mathbb R}^\perp}(\lambda|_{S_+\otimes W})^2{}^\dagger$ $(3)$, $c\,\lambda|_{S_-\otimes W^*}\lambda^*|_{S_-\otimes W}{}^\dagger$ $(3)$, $c\,\lambda|_{S_+\otimes W}\lambda^*|_{S_+\otimes W^*}{}^\dagger$ $(3)$\\
$(\lambda|_{S_-\otimes W^*},\lambda|_{S_-\otimes W^*})$ & $3$ & $(\lambda|_{S_-\otimes W^*})^2$ $(2)$, $A|_V(\lambda|_{S_-\otimes W^*})^2{}^\dagger$ $(3)$, $A|_{W_{\mathbb R}^\perp}\lambda|_{S_-\otimes W^*}\lambda|_{S_+\otimes W}{}^\dagger$ $(3)$, $c\,\lambda|_{S_-\otimes W^*}\lambda^*|_{S_+\otimes W^*}{}^\dagger$ $(3)$\\
\end{tabular}
\tabcaption{The leg sets of the special-edge composite cycles through the ten pairs of field letters at the pair weight without derivatives. The leg sets without ${}^\dagger$ enter the translation-invariant part along $\Delta$ as well.\label{tab:c11-fieldpairs}}

\begin{tabular}{l|c|c|c|c}
pair & $n=0$ & $n=1$ & $n=2$ & $n=3$\\
\hline
$(A|_V,A|_V)$ & $7\,(3)$ & $18\,(7)$ & $34\,(18)$ & $67\,(34)$\\
$(A|_V,A|_{W_{\mathbb R}^\perp})$ & $7\,(2)$ & $16\,(7)$ & $32\,(16)$ & $59\,(32)$\\
$(A|_V,\lambda|_{S_+\otimes W})$ & $9\,(4)$ & $20\,(9)$ & $42\,(20)$ & $79\,(42)$\\
$(A|_V,\lambda|_{S_-\otimes W^*})$ & $9\,(4)$ & $20\,(9)$ & $42\,(20)$ & $79\,(42)$\\
$(A|_V,c)$ & $4\,(2)$ & $8\,(4)$ & $19\,(8)$ & $35\,(19)$\\
$(A|_V,A^*|_V)$ & $15\,(6)$ & $35\,(15)$ & $73\,(35)$ & $142\,(73)$\\
$(A|_{W_{\mathbb R}^\perp},A|_{W_{\mathbb R}^\perp})$ & $2\,(1)$ & $7\,(2)$ & $16\,(7)$ & $38\,(16)$\\
$(A|_{W_{\mathbb R}^\perp},\lambda|_{S_+\otimes W})$ & $4\,(1)$ & $12\,(4)$ & $30\,(12)$ & $62\,(30)$\\
$(A|_{W_{\mathbb R}^\perp},\lambda|_{S_-\otimes W^*})$ & $4\,(1)$ & $12\,(4)$ & $30\,(12)$ & $62\,(30)$\\
$(A|_{W_{\mathbb R}^\perp},c)$ & $2\,(1)$ & $7\,(2)$ & $16\,(7)$ & $32\,(16)$\\
$(A|_{W_{\mathbb R}^\perp},A^*|_V)$ & $10\,(2)$ & $26\,(10)$ & $64\,(26)$ & $127\,(64)$\\
$(\lambda|_{S_+\otimes W},\lambda|_{S_+\otimes W})$ & $4\,(1)$ & $12\,(4)$ & $28\,(12)$ & $63\,(28)$\\
$(\lambda|_{S_+\otimes W},\lambda|_{S_-\otimes W^*})$ & $6\,(1)$ & $15\,(6)$ & $36\,(15)$ & $75\,(36)$\\
$(\lambda|_{S_+\otimes W},c)$ & $4\,(1)$ & $9\,(4)$ & $20\,(9)$ & $42\,(20)$\\
$(\lambda|_{S_+\otimes W},A^*|_V)$ & $17\,(5)$ & $42\,(17)$ & $88\,(42)$ & $173\,(88)$\\
$(\lambda|_{S_-\otimes W^*},\lambda|_{S_-\otimes W^*})$ & $4\,(1)$ & $12\,(4)$ & $28\,(12)$ & $63\,(28)$\\
$(\lambda|_{S_-\otimes W^*},c)$ & $4\,(1)$ & $9\,(4)$ & $20\,(9)$ & $42\,(20)$\\
$(\lambda|_{S_-\otimes W^*},A^*|_V)$ & $17\,(5)$ & $42\,(17)$ & $88\,(42)$ & $173\,(88)$\\
$(c,c)$ & $2\,(1)$ & $4\,(2)$ & $8\,(4)$ & $19\,(8)$\\
$(c,A^*|_V)$ & $18\,(7)$ & $34\,(18)$ & $67\,(34)$ & $121\,(67)$\\
$(A^*|_V,A^*|_V)$ & $18\,(5)$ & $52\,(18)$ & $114\,(52)$ & $235\,(114)$\\
\end{tabular}
\tabcaption{The number of leg sets of the special-edge composite cycles entering the part along $P_b$ and, in brackets, the translation-invariant part along $\Delta$, for $n$ derivatives on the two letters.\label{tab:c11-counts}}

\begin{proposition}[Output bound]
\label{prop:c11-bound}
\begin{itemize}
\item (i) A chain of Definition \ref{def:c11-sp} with outputs of total Weyl weight $w$ has superficial degree $w_{X_a}+w_{X_b}-w$ before the insertion factor of $S^{(1)}(x)$; the factor $\partial_b$ adds one, and the translation-invariant part along $\Delta$ adds none.
\item (ii) The composite pair kernel along $P_b$ is a translation-invariant local functional of Weyl weight $w_{X_a}+w_{X_b}+1$; the composite pair kernel along $\Delta$ is local of Weyl weight $w_{X_a}+w_{X_b}$, the position counted at weight $-1$, at most linear in the position. Both have the pair degree and the parity $|X_a|+|X_b|$.
\item (iii) Chains with $w\le w_{X_a}+w_{X_b}+1$ contribute to the kernel along $P_b$ and to the part along $\Delta$ linear in the position, chains with $w\le w_{X_a}+w_{X_b}$ to the translation-invariant part along $\Delta$; a chain above these bounds has a vanishing limit as $\varepsilon\to0$.
\item (iv) The number of output fields is at most $w_{X_a}+w_{X_b}+1$ plus the pair degree for both kernels; the coefficients of $\varepsilon^{-m}$ of the special-edge composite cycles have Weyl weight $w_{X_a}+w_{X_b}-2m$ along $\Delta$ and $w_{X_a}+w_{X_b}+1-2m$ along $P_b$.
\end{itemize}
The same bounds hold for the composite cycles through the pair and for the arity-one counterterm, of degree one less than the pair degree: the coefficient of $\log\varepsilon$ of the arity-one counterterm of $X_aX_b$ has Weyl weight $w_{X_a}+w_{X_b}$ along $\Delta$ and $w_{X_a}+w_{X_b}+1$ along $P_b$.
\end{proposition}

\begin{proof}
(i) Every edge obeys the weight rule $w_1+w_2-4$: $P$ (Proposition \ref{prop:c3-P}), $K_0$ (a field and its antifield have the weights $w$ and $4-w$, momentum degree $0$) and the composite kernel (Lemma \ref{lem:c11-kernels} (vi)); the count of Lemma \ref{lem:c10-degree} applies. (ii) The coefficient of $\varepsilon^0$ is a polynomial in the output momenta of degree equal to the superficial degree (Theorem \ref{thm:c11-simplex}); along $\Delta$ the position enters once, through the factor $2x^\mu$ of Lemma \ref{lem:c11-kernels} (iv). Degree and parity: Proposition \ref{prop:c11-census} (ii). (iii) A polynomial of negative degree vanishes, and so does the limit of a chain of negative degree with a special edge (Theorem \ref{thm:c11-simplex}). (iv) Weyl weight plus degree is at least $1$ for every output (Proposition \ref{prop:c10-bound}), so the number of outputs is at most $w$ plus the pair degree. The statements on the arity-one counterterm follow in the same way from Lemma \ref{lem:c4-local}.
\end{proof}

\section{Closed forms and locality}

\begin{definition}[Position part and translation-invariant part]
\label{def:c11-canonical}
A local functional is in canonical form if it is written $\int\omega\,D$ with $\omega\in\Omega^4_c(\mathbb R^4)$ and $D$ a polynomial in letters, every derivative on a field. A map $M$ of local functionals, linear in $\omega$, of the form
\[
M\Big(\int\omega\,D\Big)=\int\omega\,\big(x^\mu M_\mu(D)+M_0(D)\big),
\]
with $M_\mu$, $M_0$ maps of densities at one point and $x^\mu$ the coordinates of $V$ in an orthonormal basis, has the position part $b^\mu M_\mu$ along $b\in V$ and the translation-invariant part $M_0$. A local functional $G=\int\big(x^\mu g_\mu+g_0\big)$ at most linear in the position has the position part $\int b^\mu g_\mu$ along $b$ and the translation-invariant part $\int g_0$. On plane waves, with the momenta of the fields independent and $\omega$ absorbing the total momentum, a graph is $\alpha\,\delta(\cdot)+\beta^\mu\partial_\mu\delta(\cdot)$, and $M_0$ corresponds to $\alpha$, $M_\mu$ to $i\beta^\mu$.
\end{definition}

\begin{lemma}[Translations and the canonical form]
\label{lem:c11-canonical}
$M$ is well defined, independently of integrations by parts against $\omega$, exactly when every $M_\mu$ commutes with the derivatives and $[M_0,\partial_\mu]=M_\mu$. Then $[M,\{S^{(1)}(P_b),-\}]$ is the map $\int\omega\,D\mapsto\int\omega\,b^\mu M_\mu(D)$. For $M=v.=\{S^{(1)}(\Delta),-\}$ the position part is $2\partial_b$ and the translation-invariant part is the multiplication by $2w_D$; $\{S^{(1)}(P_b),-\}$ has position part $0$.
\end{lemma}

\begin{proof}
$M(\int\omega\,\partial_\mu D)$ and $M(-\int(\partial_\mu\omega)D)=\int\omega\big(M_\mu(D)+x^\nu\partial_\mu M_\nu(D)+\partial_\mu M_0(D)\big)$ agree for all $\omega$ exactly under the two conditions. Then $M(\int\omega\,\partial_bD)-\{S^{(1)}(P_b),M(\int\omega\,D)\}=\int\omega\,\big(M_0(\partial_bD)-\partial_bM_0(D)\big)=\int\omega\,b^\mu M_\mu(D)$, since $\{S^{(1)}(P_b),-\}$ acts by $\partial_b$ on the fields and leaves the position and $\omega$ untouched. Lemma \ref{lem:c9-SDelta} (i) gives $v.\int\omega\,D=\int\omega\,2(x\cdot\partial+w_D)D$.
\end{proof}

\begin{proposition}[Position decomposition]
\label{prop:c11-position}
\begin{itemize}
\item (i) For every pair of letters, the position part along $b$ of the composite pair kernel along $\Delta$ is twice the composite pair kernel along $P_b$. For every local functional the position part of its arity-one counterterm along $\Delta$ is twice its arity-one counterterm along $P_b$; on cocycles the position part of the composite term along $\Delta$ is twice the composite term along $P_b$; and the position part of $I^{CT}_{\mathfrak{psl}(4|4)}[\varepsilon]^{(1)}(\Delta)$ along $b$ is $2I^{CT}_{\mathfrak{psl}(4|4)}[\varepsilon]^{(1)}(P_b)$. In particular the quadratic part of the coefficient of $\log\varepsilon$ in $I^{CT}_{\mathfrak{psl}(4|4)}[\varepsilon]^{(1)}(P_b)$ is $\tfrac12$ unit times $Q\big(\int(\iota_bA|_V,c^*)-\tfrac12\int(\rho(b)\lambda^*,\lambda^*)\big)$, and the quadratic part of the arity-one obstruction along $P_b$ vanishes.
\item (ii) The translation-invariant part of the composite pair kernel along $\Delta$ is computed with the internal vertex $\int\big(2(x-x_O)^\mu\partial_\mu(A|_V)+2A|_V,A^*|_V\big)$, $x_O$ the point of the density, and pair by pair its commutator with $\{S^{(1)}(P_b),-\}$ is twice the composite pair kernel along $P_b$.
\end{itemize}
\end{proposition}

\begin{proof}
(i) By Lemma \ref{lem:c11-kernels} (iv) the internal vertex along $\Delta$ is $2\int x^\mu s_\mu+\int s_0$, with $\int s_\mu$ the internal vertex along $P_{e_\mu}$ and $s_0=2(A|_V,A^*|_V)$; insert $x^\mu=x_O^\mu+(x-x_O)^\mu$. On plane waves the factor $x^\mu$ is $-i\partial/\partial r^\mu$ of $e^{i(r,x)}$, with the momentum $r$ injected at the internal vertex and absorbed by $\omega$; its action on the momentum delta function is the position part, $-i$ times the graph with $\int s_\mu$ at $r=0$, which is $2\cdot(-i)$ times the graph along $P_{e_\mu}$; its action on the integrand is translation invariant. The same holds graph by graph for the composite cycles, their singular parts, the counterterm trees and the composite loops of $I$. Theorem \ref{thm:c9-quadratic} (i) gives the quadratic part along $\Delta$, which is $Q$ of $\int x^\mu$ times a density with $Q$ mapping letters to letters; Theorem \ref{thm:c9-quadratic} (ii) gives a quadratic obstruction along $\Delta$ without explicit position. (ii) The first statement is the decomposition of (i); the second is Lemma \ref{lem:c11-canonical} applied to the composite pair kernel along $\Delta$, whose position part is given by (i).
\end{proof}

\begin{theorem}[Special-edge composite cycles as simplex integrals]
\label{thm:c11-simplex}
For a chain of Definition \ref{def:c11-sp} with $k$ vertices of $I$ and edges $e_0,\dots,e_k$, $e_0$ at $X_a$, let $e_c$ be the composite kernel and $e_\sigma$, $\sigma\ne c$, the special edge. The cycle has $k+2$ Schwinger times: $t_j\in[\varepsilon,\infty)$ for $j\notin\{c,\sigma\}$, $t_\sigma\in[0,\varepsilon]$, and $t'_c,t''_c\in[\varepsilon,\infty)$ for the two segments of the composite kernel (Lemma \ref{lem:c11-kernels} (vi)), which meet at the internal vertex of $S^{(1)}(x)$. The edge $e_j$ carries $l+q_j$ as in Definition \ref{def:c10-cycle}, both segments carry $l+q_c$, and along $\Delta$ a momentum $r$ is injected at the internal vertex and carried by the second segment and by $e_{c+1},\dots,e_k$ to $X_b$ and $\omega$. The numerator is the product of the vertex functions, the edge factors of Definition \ref{def:c10-cycle}, the special-edge factor $-|l+q_\sigma|^2$ times the fibre factor of $K_0$, the segment factors $2(-,-)^{-1}$ times the inverse metric and $-i(l+q_c)(-,-)^{-1}$, and the insertion: $i(l+q_c)_b$ along $P_b$; $2$ for $2\int(A|_V,A^*|_V)$ and $2i(l+q_c)_\mu$ for $2\int x^\mu(\partial_\mu(A|_V),A^*|_V)$ along $\Delta$. With the simplex of dimension $k+1$ of the coordinates $\alpha=(\alpha_j)_{j\ne c}\cup(\alpha'_c,\alpha''_c)$, its centre, momentum form and $m$-th simplex integrand as in Definition \ref{def:c10-cycle} with $k+1$ in place of $k$, at generic output momenta and $T=\infty$,
\begin{equation}
\begin{aligned}
16\pi^2\big(\text{coefficient of }\varepsilon^{-m}\big)&=\frac1m\int_{\alpha_\sigma\text{ smallest}}\Big[\big(\min_{f\ne\sigma}\alpha_f\big)^m-\alpha_\sigma^m\Big]\big(m\text{-th simplex integrand}\big)\,d\alpha\qquad(m\ge1),\\
16\pi^2\big(\text{coefficient of }\varepsilon^0\big)&=\int_{\alpha_\sigma\text{ smallest}}\big(0\text{-th simplex integrand}\big)\,\log\frac{\min_{f\ne\sigma}\alpha_f}{\alpha_\sigma}\,d\alpha,
\end{aligned}
\label{eq:c11-simplex}
\end{equation}
the coefficient of $\log\varepsilon$ is $0$, and the remainder tends to $0$. Each coefficient is a polynomial in the output momenta and in $r$, of degree the superficial degree minus $2m$, and the limit vanishes for negative superficial degree. The coefficients of $\varepsilon^0$ are integrals of polynomials times logarithms of ratios of affine functions over rational subsimplices; for $k=1$ they lie in $\mathbb Q+\mathbb Q\log2+\mathbb Q\log3$ (Proposition \ref{prop:c11-k1}).
\end{theorem}

\begin{proof}
The domain $\{t_\sigma\le\varepsilon,\ t_f\ge\varepsilon\ (f\ne\sigma)\}$ is the difference of the domains with the cutoff shapes $1/\min_{f\ne\sigma}\alpha_f$ and $1/\min_f\alpha_f$, which differ on the region where $\alpha_\sigma$ is the smallest coordinate. Proposition \ref{prop:c10-shape} applies to each with $k+2$ edges, the Schwinger measure being $\tau^{k+1}d\tau\,d\alpha$, and the difference is taken as in Lemma \ref{lem:c10-special}: the two terms in $\log\varepsilon$ cancel, and the finite parts differ by $-\int F\log(\alpha_\sigma/\min_{f\ne\sigma}\alpha_f)$ over that region, $F$ the $0$-th simplex integrand. The region $\tau\ge1$ converges at $T=\infty$: the composite kernel is of order $|p|^{-3}$ at small momentum (Proposition \ref{prop:c9-composite} (iv)) and the other edges carry distinct generic momenta (proof of Theorem \ref{thm:c10-decomp} (ii)), so the integrand is integrable and the measure of $\{\alpha_\sigma\le\varepsilon/\tau\}$ is of order $\varepsilon$; the injected momentum enters polynomially at each order in $r$. The degree statement is Proposition \ref{prop:c11-bound} (i).
\end{proof}

\begin{proposition}[Generating functions in the point splitting]
\label{prop:c11-genfun}
With the letters at $x+y_a$ and $x+y_b$, $y=y_b-y_a$, the $m$-th simplex integrand for the numerator $(\text{chain numerator})\,e^{i(y,l)}$ is the chain numerator with $l$ replaced by $-i\partial_y$, applied to
\[
(-1)^{m+k}e^{i(y,\text{centre})}\sum_{i\ge0}\frac{(|y|^2/4)^{i+m+k}(\text{momentum form})^i}{i!\,(i+m+k)!}=(-1)^{m+k}e^{i(y,\text{centre})}\Big(\frac{|y|}{2\sqrt{\text{momentum form}}}\Big)^{m+k}I_{m+k}\big(|y|\sqrt{\text{momentum form}}\big),
\]
and, summed over the chains with the signs of Definition \ref{def:c11-sp},
\[
\sum_{\alpha,\beta}\frac{y_a^\alpha y_b^\beta}{\alpha!\,\beta!}\big(\text{composite pair kernel of }(\partial^\alpha\varphi_a,\partial^\beta\varphi_b)\text{ along }x\big)=\frac{\hbar}{16\pi^2}\,e^{i(y_b,q_k)}\int_{\alpha_\sigma\text{ smallest}}\big(0\text{-th simplex integrand}\big)\log\frac{\min_{f\ne\sigma}\alpha_f}{\alpha_\sigma}\,d\alpha,
\]
with $q_k$ including $r$ along $\Delta$; the coefficients of $\varepsilon^{-m}$ have the same generating function with the weights of \eqref{eq:c11-simplex}. In position space $iq_j$ is the derivative of the product of the outputs at $x$.
\end{proposition}

\begin{proof}
As Proposition \ref{prop:c10-genfun}, with $k+2$ edges: the power $\tau^{-1-m}$ of the Gaussian-integrated integrand times $\tau^{k+1}$ fixes the index shift $m+k$, and Theorem \ref{thm:c11-simplex} replaces the weight $(\min_e\alpha_e)^m$ by the weights of the region where $\alpha_\sigma$ is the smallest.
\end{proof}

\begin{proposition}[One vertex in closed form]
\label{prop:c11-k1}
Let $k=1$, $\beta:=\alpha_\sigma$, and let $\gamma'$ and $\gamma''=1-\beta-\gamma'$ be the coordinates of the two segments. At $r=0$ the centre $-\beta q_\sigma-(1-\beta)q_c$ and the momentum form $\beta(1-\beta)|q_c-q_\sigma|^2$ depend on $\beta$ alone, the region where $\beta$ is the smallest coordinate requires $\beta\le\tfrac13$, and
\[
\begin{aligned}
16\pi^2\big(\text{composite pair kernel}\big)&=\int_0^{1/3}\big(0\text{-th simplex integrand}\big)(\beta)\Big[(1-\beta)\log\frac{1-\beta}{2\beta}-(1-3\beta)\Big]d\beta,\\
16\pi^2\big(\text{coefficient of }\varepsilon^{-m}\big)&=\int_0^{1/3}\big(m\text{-th simplex integrand}\big)(\beta)\,\frac2m\Big[\frac{\big(\tfrac{1-\beta}2\big)^{m+1}-\beta^{m+1}}{m+1}-\beta^m\Big(\frac{1-\beta}2-\beta\Big)\Big]d\beta .
\end{aligned}
\]
The part linear in $r$ carries one factor $\gamma''$ in the simplex integrand, and $\int\gamma''(\cdots)\,d\gamma'=\tfrac{1-\beta}2\int(\cdots)\,d\gamma'$ by the symmetry $\gamma'\leftrightarrow\gamma''$ of the region and of the weights. The coefficients are integrals of $\beta^a(1-\beta)^b$ times these weights over $[0,\tfrac13]$, in $\mathbb Q+\mathbb Q\log2+\mathbb Q\log3$.
\end{proposition}

\begin{proof}
On the region, $\min_{f\ne\sigma}\alpha_f=\min(\gamma',\gamma'')$ and $\gamma'\in[\beta,1-2\beta]$. By the symmetry about $\gamma'=\tfrac{1-\beta}2$,
\[
\int\log\frac{\min(\gamma',\gamma'')}\beta\,d\gamma'=2\int_\beta^{(1-\beta)/2}\log\frac{\gamma'}\beta\,d\gamma'=(1-\beta)\log\frac{1-\beta}{2\beta}-(1-3\beta),
\]
and $\frac1m\int\big(\min(\gamma',\gamma'')^m-\beta^m\big)d\gamma'$ is the second weight. At $r=0$ both segments carry $l+q_c$, so the Gaussian data depend on $\beta$ alone; along $\Delta$ the derivative in $r$ of $\sum_jt_j|l+q_j|^2$ is $2t''_c(l+q_c)\cdot$, linear in $\gamma''$. The integrals $\int_0^{1/3}\beta^a(1-\beta)^b\log\frac{1-\beta}{2\beta}\,d\beta$ are rational combinations of $1$, $\log2$, $\log3$, the endpoint $\beta=\tfrac13$ giving $\log\tfrac23$.
\end{proof}

\begin{theorem}[The composite term is two-body and local]
\label{thm:c11-local}
Let $O=\int\omega\,X_1\cdots X_n$ be a local functional with $sO=0$, $x\in\{\Delta,P_b\}$, $T=\infty$, generic momenta, and $\hbar F_1^{(0)}$ as in (i) of Definition \ref{def:c10-hyp}.
\begin{itemize}
\item (i) The special-edge composite cycles of $(X_a,X_b)$ along $x$ are $\sum_{m\ge1}\varepsilon^{-m}$ times local coefficients, plus the composite pair kernel, plus a remainder tending to $0$; they contain no $\log\varepsilon$.
\item (ii) $s$ annihilates the coefficient of $\log\varepsilon$ in $I^{CT}_{\mathfrak{psl}(4|4)}[\varepsilon]^{(1)}(x)$, and $I^{CT}_{\mathfrak{psl}(4|4)}[\varepsilon]^{(1)}(x)$ is $\log\varepsilon$ times that coefficient in every sector.
\item (iii) (Arity-one BRST covariance.) $s$ of the coefficient of $\log\varepsilon$ in the arity-one counterterm of $O$ along $x$ is minus the bracket of $O$ with the coefficient of $\log\varepsilon$ in $I^{CT}_{\mathfrak{psl}(4|4)}[\varepsilon]^{(1)}(x)$, and the sum over the pairs, with the Koszul signs of \eqref{eq:c10-decomp}, of the coefficients of $\varepsilon^{-m}$ of the special-edge composite cycles is $s$ of the coefficient of $\varepsilon^{-m}$ in the arity-one counterterm of $O$.
\item (iv) The composite term of $O$ along $\Delta$ (Definition \ref{def:c10-hyp}) and along $P_b$ (Proposition \ref{prop:c10-P1}) is
\begin{equation}
\big(\text{composite term of }O\text{ along }x\big)=\sum_{a<b}(-1)^{|X_b|(|X_{a+1}|+\cdots+|X_{b-1}|)}\int\omega\,X_1\cdots X_{a-1}\,\big(\text{composite pair kernel of }(X_a,X_b)\text{ along }x\big)\,X_{a+1}\cdots\widehat{X_b}\cdots X_n,
\label{eq:c11-local}
\end{equation}
a local two-body functional with no one-body term.
\item (v) Hence the composite terms of the pairs and of the legs of Theorem \ref{thm:c10-D1}, summed, equal the sum over the pairs of the composite pair kernels along $\Delta$: the nonlocal parts of the brackets of the letters $A|_V$, $\lambda$, $c$ with the renormalised composite two-point functional, $\log|p|^2$ of the leg momentum and $(p,a)(p,b)|p|^{-2}$ on $A|_V$ (Remark \ref{rem:c9-Pi}), cancel against those of the images under $Q$ of the renormalised composite cycles through the pairs containing these letters.
\item (vi) For $sO\ne0$ the right side of \eqref{eq:c11-local} acquires the renormalised composite graphs of $sO$, minus $\tfrac\hbar2\,\partial$ of the composite kernel applied to the trees with one vertex $sO$; these contain the nonlocal bracket with the renormalised composite two-point functional on the letters of $sO$.
\end{itemize}
\end{theorem}

\begin{proof}[Outline]
By Proposition \ref{prop:c11-master}, $Q$ of the composite graphs of $O$ is $\tfrac\hbar2\,\partial$ of the composite kernel applied to the trees with one edge $K_\varepsilon-K_0$. Sort them by the position of the special edge relative to the cycle closed by the composite kernel: (A) on a cycle through two letters, the special-edge composite cycles with trees on the outputs; (B) on a tree attached to a cycle through two letters, at an output or at a spectator letter, the factor $e^{-\varepsilon|p|^2}-1$ at a generic momentum times the unrenormalised composite cycle; (C) on a composite loop of $I$ attached to $O$ by a tree: by Proposition \ref{prop:c11-master} for $M=0$, read on the amputated loops, where the collapse of a tree edge at a leg of the loop by $QP\star+P\star Q=e^{-\varepsilon H}$ gives the bracket with $I$, minus the amputated loops with one special edge is $s$ of the amputated composite loops of $I$, that is, $s$ of the renormalised composite functional along $x$ plus $s\,I^{CT}_{\mathfrak{psl}(4|4)}[\varepsilon]^{(1)}(x)$, and (C) is its tree insertion into $O$; (D) on the tree joining a composite loop of $I$ to $O$, the factor $e^{-\varepsilon|p|^2}-1$ times the unrenormalised loop. The counterterm trees, with the vertices $I^{CT}_{\mathfrak{psl}(4|4)}[\varepsilon]^{(1)}(x)$ and the arity-one counterterm of $O$, have by Proposition \ref{prop:c4-defectflow} and Lemma \ref{lem:c6-bvexact} the images under $Q$
\[
\begin{aligned}
Q\big(-\text{tree}(I^{CT,(1)},O)\big)&=-\text{tree}\big(sI^{CT,(1)},O\big)-\{I^{CT,(1)},O\}+(\text{special edges}),\\
Q\big(-\text{tree}(\text{arity-one counterterm of }O)\big)&=-\text{tree}\big(s(\text{arity-one counterterm of }O)\big)+(\text{special edges}),
\end{aligned}
\]
with $I^{CT,(1)}$ short for $I^{CT}_{\mathfrak{psl}(4|4)}[\varepsilon]^{(1)}(x)$, an odd vertex, and the special edges joining the counterterm vertex to a vertex $I$. These special edges cancel the singular parts of (B) and (D), whose sums are $e^{-\varepsilon|p|^2}-1$ times a renormalised cycle or loop, which have limits, and tend to $0$.

(ii) By (C), $s$ of the amputated composite loops of $I$ is minus the loops with one special edge, which contain no $\log\varepsilon$ (Theorem \ref{thm:c11-simplex}). The coefficient of $\varepsilon^{-m}$ of $I^{CT}_{\mathfrak{psl}(4|4)}[\varepsilon]^{(1)}(x)$ is a colour-singlet local functional of degree $-1$ and Weyl weight $4-2m$ along $\Delta$, the position at weight $-1$ and at most linear in it, and $5-2m$ along $P_b$; a colour singlet has at least two adjoint fields, and degree $-1$ requires an antifield of degree $-1$, of weight at least $\tfrac52$, and a field of degree $0$, of weight at least $1$, or $c^*$ and $c$, of weights $4$ and $0$: the field weight is at least $\tfrac72$, and at least $\tfrac52$ with one position, above $4-2m\le2$ and $5-2m\le3$.

(iii), (iv) The loops (C) give the tree insertion of $s$ of the renormalised composite functional, and by (ii) $-\text{tree}(sI^{CT,(1)},O)$ vanishes; $s$ of the renormalised composite functional is the composite part (Remark \ref{rem:c9-Pi}, Proposition \ref{prop:c9-structure} (ii)), whose tree insertion is subtracted in the definition of the composite term. What remains is (A) minus the trees of $s$ of the arity-one counterterm of $O$ minus $\{I^{CT,(1)},O\}$. The left side has a limit, the composite graphs minus their counterterm trees converging (Proposition \ref{prop:c9-structure} (i)), and (A) contains no $\log\varepsilon$ by (i): the coefficient of $\log\varepsilon$ gives (iii), and so do the coefficients of $\varepsilon^{-m}$. At order $\varepsilon^0$ the trees on the outputs of the special-edge composite cycles and of the poles of the arity-one counterterm are the same trees with edges $P(\varepsilon,\infty)$, so after the cancellation of the poles their difference tends to the trees of the composite pair kernels with edges $P(0,\infty)$, and the injectivity of the tree-level flow gives \eqref{eq:c11-local}. Statement (i) is Theorem \ref{thm:c11-simplex} for each chain.

(v) Definition \ref{def:c10-hyp} writes the same limit as the sum of the composite terms of the pairs and of the legs (proof of Theorem \ref{thm:c10-D1}), so the two sums agree on cocycles, and the right side is local. The mechanism: $Q$ of a renormalised composite cycle through $(X_a,X_b)$ produces, through $QP\star+P\star Q=e^{-\varepsilon H}$ on the edge at $X_a$, the composite cycles of the letters of $\{I,X_a\}$, and $sO=0$ trades them for the composite loops on the legs, whose renormalised value is the renormalised composite two-point functional. (vi) For $sO\ne0$, Proposition \ref{prop:c11-master} keeps the trees with one vertex $sO$, whose composite closures are the composite graphs of $sO$ of both kinds.
\end{proof}

\begin{remark}[What the locality uses and what it fixes]
\label{rem:c11-local}
The argument uses $sO=0$ exactly, the classical master equation, the anticommutation of $Q$ with $\partial$ of the composite kernel, which rests on $[v_x.,Q]=0$ and $QP\star+P\star Q=e^{-\varepsilon H}$ (Lemma \ref{lem:c11-kernels} (v)), the absence of $\log\varepsilon$ in cycles with one special edge, a property of the cutoff on every edge (Theorem \ref{thm:c11-simplex}), the absence of a term $\varepsilon^{-1}$ in the composite loops of $I$ (Theorem \ref{thm:c11-local} (ii)), and generic momenta at $T=\infty$ for the convergence of the cycles and of the composite loops, which are of order $|p|^{-3}$ (Proposition \ref{prop:c9-composite} (iv)), or any finite $T$ (Theorem \ref{thm:c8-L} (ii)); (i) of Definition \ref{def:c10-hyp} enters through the terms with $QW$ of the proof of Theorem \ref{thm:c10-D1}, which are outside the composite term. For a functional of the fields $A|_V$, $A|_{W_{\mathbb R}^\perp}$, $\lambda$ alone, $sO=0$ is the gauge invariance of its density. Statement (iv) fixes the one-body part of the composite term: a term $t$ units times $v.X_a$ on every letter, added to the right side of \eqref{eq:c11-local}, would replace each sector coefficient $\theta$ by $\theta+t$ in Sections 11.3--11.5 and leave every spectrum unchanged; the value $\theta=-\tfrac13$ of Proposition \ref{prop:c11-theta-sector} is the value in the normalisation of \eqref{eq:c11-local}.
\end{remark}

\begin{corollary}[Local two-body form of $D_1$ and $P_1$]
\label{cor:c11-D1local}
Under the hypotheses of Definition \ref{def:c10-hyp}, and of the same definition with $P_b$ in place of $\Delta$, for $O=\int\omega\,X_1\cdots X_n$, with the Koszul signs of \eqref{eq:c10-decomp},
\[
\begin{aligned}
D_1(O)=\ &\sum_{a<b}\pm\int\omega\cdots\Big[4\big(\text{coefficient of }\log\varepsilon\text{ in }\partial_\delta I^{CT}[\varepsilon](X_aX_b)\big)+\big(\text{composite pair kernel of }(X_a,X_b)\text{ along }\Delta\big)\Big]\cdots\\
&+\sum_a\int\omega\cdots\Big[-4\,\text{units}\cdot(\text{weight of }X_a)\,X_a-\big\{\hbar F_1^{(1)}(\Delta)-4\,\text{units}\cdot(\text{field-counting functional}),X_a\big\}\Big]\cdots,\\
P_1(O)=\ &\sum_{a<b}\pm\int\omega\cdots\big(\text{composite pair kernel of }(X_a,X_b)\text{ along }P_b\big)\cdots-\sum_a\int\omega\cdots\big\{\hbar F_1^{(1)}(P_b),X_a\big\}\cdots,
\end{aligned}
\]
every term local; $D_1$ minus four times the map of Theorem \ref{thm:c10-brst} is the sum over the pairs of the composite pair kernels along $\Delta$ minus the bracket of $O$ with $\hbar F_1^{(1)}(\Delta)-4$ units times the field-counting functional.
\end{corollary}

\begin{proof}
Insert \eqref{eq:c11-local} into Theorem \ref{thm:c10-D1} and Proposition \ref{prop:c10-P1}; four times the coefficient of $\log\varepsilon$ is two-body by Theorem \ref{thm:c10-decomp} (iii), and the brackets with the field-counting functional, with $\hbar F_1^{(1)}(\Delta)$ and with $\hbar F_1^{(1)}(P_b)$ are derivations.
\end{proof}

\begin{proposition}[The route of the special edges]
\label{prop:c11-routes}
For $O$ as in Corollary \ref{cor:c11-D1local}, the limit as $\varepsilon\to0$ of the cycles through the pairs of letters of $O$ with one special edge and one vertex $S^{(1)}(\Delta)$ anywhere on the cycle, as the composite kernel or with one leg on the special edge, minus the trees of $s$ of the arity-one counterterm of $O$, the graphs of Lemma \ref{lem:c8-transfer} used in Theorem \ref{thm:c9-D1}, is the sum over the pairs of four times the coefficient of $\log\varepsilon$ of the pair kernel plus the composite pair kernel along $\Delta$.
\end{proposition}

\begin{proof}
Both routes compute the local representative $D_1(O)$, which is unique (Theorem \ref{thm:c8-L} (i)). In the route of the special edges the one-body terms are the trees with $-\{\hbar F_1^{(1)}(\Delta),O\}$, the arity-one defect of the action joined to $O$ vanishing by \eqref{eq:c9-arity1} (proof of Theorem \ref{thm:c9-D1}); in Corollary \ref{cor:c11-D1local} the one-body terms are $-4$ units times the weights minus the bracket with $\hbar F_1^{(1)}(\Delta)-4$ units times the field-counting functional, which is $-\{\hbar F_1^{(1)}(\Delta),-\}$. The two-body sums therefore agree.
\end{proof}

\begin{example}[The scalar bilinears]
\label{ex:c11-bilinear}
For $O=\int\omega\,(\iota_{w_1}A,\iota_{w_2}A)$, on the part with two legs $A|_{W_{\mathbb R}^\perp}$, Theorem \ref{thm:c9-D1} and Remark \ref{rem:c9-twoeval} give the value $4\,(\text{coefficient of }\log\varepsilon\text{ in }\partial_\delta I^{CT}[\varepsilon](O))-\tfrac43\,O-\tfrac13\{S^{(1)}(\Delta),O\}$ units of the route of the special edges, so by Proposition \ref{prop:c11-routes}
\[
\big(\text{composite pair kernel of }(A|_{W_{\mathbb R}^\perp},A|_{W_{\mathbb R}^\perp})\text{ along }\Delta\big)\Big|_{\text{outputs }A|_{W_{\mathbb R}^\perp}A|_{W_{\mathbb R}^\perp}}=\Big[-\tfrac13\,v.O-\tfrac43\,O\Big]\ \text{units},
\]
for the traceless and the trace parts alike. It is the single chain of Table \ref{tab:c11-fieldpairs},
\[
A|_{W_{\mathbb R}^\perp}\xrightarrow{\ K_\varepsilon-K_0\ }\big([A|_{W_{\mathbb R}^\perp},c],A^*|_{W_{\mathbb R}^\perp}\big)\xrightarrow{\ \text{composite kernel}\ }A|_V(A|_{W_{\mathbb R}^\perp})^2\xrightarrow{\ P\ }A|_{W_{\mathbb R}^\perp}.
\]
Its position part along $b$ is $-\tfrac23$ units times $\partial_b$, so the composite pair kernel along $P_b$ is $-\tfrac13$ units times $\{S^{(1)}(P_b),O\}$, in agreement with Theorem \ref{thm:c9-P1}, the composite triangle and the composite two-point functional along $P_b$ giving $-2+\tfrac53$; its translation-invariant part is $-\tfrac23\cdot2\,O-\tfrac43\,O=-\tfrac83\,O$ units.
\end{example}

\section{The free directions and the arity-two equation}

\begin{definition}[Finite counterterms along $\Delta$ and $P_b$]
\label{def:c11-free}
Write the solutions of the arity-one equations along $\Delta$ and along $P_b$ as
\[
\begin{aligned}
\hbar F_1^{(1)}(\Delta)=\ &4\,\text{units}\cdot(\text{field-counting functional})+\Big[-\tfrac23\int(A,A^*)+\tfrac13\int(\lambda^*,\lambda)+\tfrac43\int(c^*,c)\Big]\,\text{units}\\
&+\theta\,\text{units}\int\big(2E(A|_{W_{\mathbb R}^\perp})+2A|_{W_{\mathbb R}^\perp},A^*|_{W_{\mathbb R}^\perp}\big)-\theta_\lambda\,\text{units}\int\big(2E(\lambda)+3\lambda,\lambda^*\big)\\
&+\theta_V\,\text{units}\Big[\int\big(2E(A|_V)+2A|_V,A^*|_V\big)-\int\big(2E(c),c^*\big)\Big]+\text{units}\cdot Q\Big(\kappa_V\int(\iota_xA|_V,c^*)+\kappa_\lambda\int(\rho(x)\lambda^*,\lambda^*)\Big)+(\text{cubic and higher}),\\
\hbar F_1^{(1)}(P_b)=\ &\theta_P\,\text{units}\int\big(\partial_b(A|_{W_{\mathbb R}^\perp}),A^*|_{W_{\mathbb R}^\perp}\big)-\theta_{P,\lambda}\,\text{units}\int(\partial_b\lambda,\lambda^*)+\theta_{P,V}\,\text{units}\Big[\int\big(\partial_b(A|_V),A^*|_V\big)-\int(\partial_bc,c^*)\Big]\\
&+Q\big(\text{a combination of }\textstyle\int(\iota_bA|_V,c^*)\text{ and }\int(\rho(b)\lambda^*,\lambda^*)\big)+(\text{cubic and higher}),
\end{aligned}
\]
with numbers $\theta,\theta_\lambda,\theta_V,\kappa_V,\kappa_\lambda,\theta_P,\theta_{P,\lambda},\theta_{P,V}$; $\theta$ is the coefficient of \eqref{eq:c9-Fperp} and $\theta_P$ that of Theorem \ref{thm:c9-P1}. The three terms with $\theta$, $\theta_\lambda$, $\theta_V$ are the sector parts of $S^{(1)}(\Delta)$ in the sectors of $A|_{W_{\mathbb R}^\perp}$, of $\lambda$ and of $(A|_V,c)$ times these numbers, and those with $\theta_P$, $\theta_{P,\lambda}$, $\theta_{P,V}$ the sector parts of $S^{(1)}(P_b)$ (Definition \ref{def:c11-kernels}). The scheme is translation covariant if the position part along $b$ of $\hbar F_1^{(1)}(\Delta)$ is $2\hbar F_1^{(1)}(P_b)$; then $\theta_P=\theta$, $\theta_{P,\lambda}=\theta_\lambda$, $\theta_{P,V}=\theta_V$.
\end{definition}

\begin{proposition}[Completeness of the quadratic directions]
\label{prop:c11-tc0}
Let $\mathcal F_{-1}$ and $\mathcal F_{-2}$ be the quadratic local functionals of degree $-1$ and $-2$, of Weyl weight $4$ with the position at weight $-1$, at most linear in the position and invariant under $\mathrm{Spin}(4)\times\mathrm{Spin}(6)$. On the subspaces invariant under the reflection of Theorem \ref{thm:c9-quadratic} (iii),
\[
\frac{\ker\big(Q|_{\mathcal F_{-1}}\big)}{Q\,\mathcal F_{-2}}=\operatorname{span}\big\{\text{the three sector parts of }S^{(1)}(\Delta)\big\},\qquad Q\,\mathcal F_{-2}=\operatorname{span}\Big\{Q\int(\iota_xA|_V,c^*),\ Q\int(\rho(x)\lambda^*,\lambda^*)\Big\}.
\]
Without the reflection the quotient has one further class, $\int(\lambda^*,\rho(e_1e_2e_3e_4)\lambda)$, with $e_1,\dots,e_4$ an oriented orthonormal basis of $V$.
\end{proposition}

\begin{proof}
By Schur's lemma, $\mathcal F_{-1}$ is spanned by $\int(A^*|_V,D\,A|_V)$ with the structures $D$ built from $1$, $x\cdot\partial$, $x_\mu\partial_\nu$, $x_\nu\partial_\mu$ on the vector index and $\epsilon_{\mu\nu\kappa\tau}x^\kappa\partial^\tau$; by $\int(A^*|_{W_{\mathbb R}^\perp},D\,A|_{W_{\mathbb R}^\perp})$ with $D\in\{1,x\cdot\partial\}$; by $\int(\lambda^*,D\lambda)$ with $D\in\{1,\ x\cdot\partial,\ \rho(x)\rho(\partial)\}$ times $\{1,\rho(e_1e_2e_3e_4)\}$; and by $\int(c^*,Dc)$ with $D\in\{1,x\cdot\partial\}$. $\mathcal F_{-2}$ is spanned by $\int x^\mu(A_\mu,c^*)$, $\int(\lambda^*,\rho(x)\lambda^*)$ and $\int(\lambda^*,\rho(e_1e_2e_3e_4)\rho(x)\lambda^*)$. The structures with $\epsilon$ or with one factor $\rho(e_1e_2e_3e_4)$ change sign under the reflection. With $Q$ of Proposition \ref{prop:c1-Q}, the ranks on the plane-wave normal forms, computed exactly, are, in the sectors of $A|_{W_{\mathbb R}^\perp}$, of $\lambda$ and of $(A|_V,c)$: $\dim\mathcal F_{-1}=2,3,6$, $\dim\ker Q=1,2,2$, $\dim Q\mathcal F_{-2}=0,1,1$ for the reflection-invariant subspaces. The kernel is spanned by the sector part of $A|_{W_{\mathbb R}^\perp}$; by the sector part of $\lambda$ and $Q\int(\rho(x)\lambda^*,\lambda^*)$; by the sector part of $(A|_V,c)$ and $Q\int(\iota_xA|_V,c^*)$. Without the reflection, $\dim\mathcal F_{-1}=15$, the $\epsilon$-structure is not $Q$-closed, $\int(\lambda^*,\rho(e_1e_2e_3e_4)\lambda)$ is closed and not exact, and $\int(\lambda^*,\rho(e_1e_2e_3e_4)\rho(x)\lambda^*)=0$.
\end{proof}

\begin{proposition}[What the arity-one equations fix]
\label{prop:c11-arity1}
\begin{itemize}
\item (i) Along $\Delta$ the quadratic part of \eqref{eq:c9-arity1} fixes the translation-invariant part $\big[-\tfrac23\int(A,A^*)+\tfrac13\int(\lambda^*,\lambda)+\tfrac43\int(c^*,c)\big]$ units (Theorem \ref{thm:c9-quadratic} (iii)) and leaves $\theta$, $\theta_\lambda$, $\theta_V$, $\kappa_V$, $\kappa_\lambda$ free. Along $P_b$ the quadratic part of the arity-one obstruction vanishes (Proposition \ref{prop:c11-position} (i)), so the quadratic part of $\hbar F_1^{(1)}(P_b)$ is $Q$-closed and every coefficient of Definition \ref{def:c11-free} along $P_b$ is free; the translation-covariant choice, $\hbar F_1^{(1)}(P_b)$ equal to half the position part along $b$ of $\hbar F_1^{(1)}(\Delta)$, solves the equation along $P_b$ at every order in the fields.
\item (ii) Two solutions along $\Delta$ differ by an $s$-closed local functional of degree $-1$ and Weyl weight $4$. The common shift of $\theta$, $\theta_\lambda$, $\theta_V$ by $t$ adds $t$ units times $S^{(1)}(\Delta)$, which is $s$-closed (Lemma \ref{lem:c9-SDelta} (iii)); the directions $Q\int(\iota_xA|_V,c^*)$ and $Q\int(\rho(x)\lambda^*,\lambda^*)$ complete to $s\int(\iota_xA|_V,c^*)$ and $s\int(\rho(x)\lambda^*,\lambda^*)$, and an $s$-exact direction $sG$ contributes $s\{G,O\}$ to $D_1(O)$ on cocycles, zero on cohomology.
\item (iii) Since $\{I,S^{(1)}(\Delta)\}=0$ while the brackets of $I$ with the sector parts are nonzero, the cubic part of \eqref{eq:c9-arity1} depends on the sector coefficients through $(\theta-\theta_V)$ times the bracket of the cubic part of $I$ with the sector part of $A|_{W_{\mathbb R}^\perp}$ plus $(\theta_\lambda-\theta_V)$ times its bracket with the sector part of $\lambda$. It is solvable exactly when this sum is congruent to the cubic part of the composite part minus the bracket of the cubic part of $I$ with the translation-invariant part, modulo the images under $Q$ of the cubic local functionals of degree $-1$. If the classes of the two brackets modulo the images under $Q$ of the cubic local functionals of degree $-1$, Weyl weight $4$ and degree at most one in the position are linearly independent, the cubic equation fixes $\theta-\theta_V$ and $\theta_\lambda-\theta_V$.
\end{itemize}
\end{proposition}

\begin{proof}
(i) Theorem \ref{thm:c9-quadratic} (iii) along $\Delta$, Proposition \ref{prop:c11-position} (i) along $P_b$. For the covariant choice: $s$ maps letters to letters and carries no position, so it commutes with taking the position part, and the position part of the arity-one obstruction along $\Delta$ is twice the obstruction along $P_b$ by Proposition \ref{prop:c11-position} (i), the term $4$ units times $s$ of the field-counting functional having no position part. (ii) Theorem \ref{thm:c8-obstruction} (iii) restricted to arity one along $\Delta$, and Lemma \ref{lem:c10-deriv}: $\{sG,O\}=-s\{G,O\}$ for $sO=0$ and $G$ of degree $-2$. (iii) $s$ of the sector terms is the sum of $\theta$, $\theta_\lambda$, $\theta_V$ units times the brackets of $I$ with the sector parts, whose sum over the three sectors vanishes; at cubic order $s$ of the solution is $Q$ of its cubic part plus the bracket of the cubic part of $I$ with its quadratic part.
\end{proof}

\begin{theorem}[The position part of $D_1$ is exact]
\label{thm:c11-arity2}
Assume a one-loop quantisation through arity two along $\Delta$ and $P_b$, that is, \eqref{eq:c8-qme} modulo $\hbar^2$ in its components of arity at most two along these generators, and local representatives through arity two along $(\Delta,P_b)$ on a space $\mathfrak F$ of local functionals stable under $\Lambda$, with $P_1$ translation covariant, $P_1(\int\omega\,D)=\int\omega\,P_1(D)$ for a map $P_1$ of densities commuting with the derivatives; this holds when $\hbar F_1^{(1)}(P_b)$ is translation invariant. Then for every $O=\int\omega\,D\in\mathfrak F$ with $sO=0$
\[
\big(\text{position part of }D_1\text{ along }b\big)(O)=\big[D_1,\{S^{(1)}(P_b),-\}\big](O)=\pm\,s\Big(\hbar L_1(O)^{(2)}(\Delta,P_b)\Big):
\]
the position part of $D_1$ is $s$-exact on cocycles, with the arity-two component of $O$ as primitive.
\end{theorem}

\begin{proof}
By Corollary \ref{cor:c8-homotopy} (ii), $[\Lambda^{(1)}(\Delta),\Lambda^{(1)}(P_b)]+\Lambda^{(1)}([\Delta,P_b])$ is a graded commutator of $\Lambda^{(0)}$ with an operator built from $\Lambda^{(2)}(\Delta,P_b)$. With $[\Delta,P_b]=-2P_b$ (Remark \ref{rem:c2-dilatation}) and $S^{(2)}(\Delta,P_b)=0$ (Definition \ref{def:c2-S2}), the order-$\hbar^0$ part of that operator vanishes and the order-$\hbar$ part of the identity reads
\[
\big[D_1,\{S^{(1)}(P_b),-\}\big]+\big[v.,P_1\big]-2P_1=\big[s,H_1\big],
\]
with $H_1$ built from $\hbar L_1^{(2)}(\Delta,P_b)$. For translation-covariant $P_1$, $v.(P_1O)=\int\omega\,2(x\cdot\partial+w_D+1)P_1(D)$ and $P_1(v.O)=\int\omega\,\big(2x^\mu P_1(\partial_\mu D)+2w_DP_1(D)\big)$, so $[v.,P_1]=2P_1$. Lemma \ref{lem:c11-canonical} identifies $[D_1,\{S^{(1)}(P_b),-\}]$ with the position part, and on a cocycle $[s,H_1]O=s(H_1O)$.
\end{proof}

\begin{lemma}[Position part of a bracket]
\label{lem:c11-bracketpos}
For a local functional $G=\int(x^\mu g_\mu+g_0)$, the position part along $b$ of $\{G,-\}$ is $\{\int b^\mu g_\mu,-\}$.
\end{lemma}

\begin{proof}
The bracket contracts a slot of $g_\mu$ at $y$ with a letter $\partial^\alpha\varphi$ of $D$ at $x$ through $K_0$, which gives $\partial^\alpha_x(x^\mu\cdots)$; the term with no derivative on $x^\mu$ is $x^\mu$ times the bracket with $\int g_\mu$, and the others carry no position.
\end{proof}

\begin{corollary}[The composite term along $P_b$ and the translation component]
\label{cor:c11-theta}
Under the hypotheses of Theorem \ref{thm:c11-arity2} and of Definition \ref{def:c10-hyp}, on every cocycle $O\in\mathfrak F$:
\begin{itemize}
\item (i) twice the composite term of $O$ along $P_b$ minus the bracket of $O$ with the position part along $b$ of $\hbar F_1^{(1)}(\Delta)$ is $s$-exact; this position part is
\[
2\theta\,\text{units}\int\big(\partial_b(A|_{W_{\mathbb R}^\perp}),A^*|_{W_{\mathbb R}^\perp}\big)-2\theta_\lambda\,\text{units}\int(\partial_b\lambda,\lambda^*)+2\theta_V\,\text{units}\Big[\int\big(\partial_b(A|_V),A^*|_V\big)-\int(\partial_bc,c^*)\Big]+\text{units}\cdot Q\Big(\kappa_V\int(\iota_bA|_V,c^*)+\kappa_\lambda\int(\rho(b)\lambda^*,\lambda^*)\Big)
\]
plus the position parts of the higher terms;
\item (ii) if $\theta=\theta_\lambda=\theta_V$, the composite term along $P_b$ induces $\theta$ units times $\{S^{(1)}(P_b),-\}$ on the cohomology of $\mathfrak F$, and on every family in which $\{S^{(1)}(P_b),O\}$ is not exact for some cocycle the composite term along $P_b$ determines $\theta$;
\item (iii) in the translation-covariant scheme the position part of $D_1$ is $2P_1$, and the arity-two equation along $(\Delta,P_b)$ is equivalent to $P_1\cong0$ on cohomology: $P_1(O)=\pm\tfrac12\,s\big(\hbar L_1(O)^{(2)}(\Delta,P_b)\big)$;
\item (iv) if $\theta_P=\theta_{P,\lambda}=\theta_{P,V}$, then $P_1$ induces $(\theta-\theta_P)$ units times $\{S^{(1)}(P_b),-\}$ on cohomology, and the finite part $\hbar\psi'=\tfrac a2\{S^{(1)}(\Delta),-\}$, $a=(\theta-\theta_P)$ units, of Proposition \ref{prop:c9-scheme} gives $P_1\cong0$ on every family at once, leaves $D_1$, $M_1$, $R_1$ unchanged and changes the components along $K_b$, $\Sigma$ and $\Sigma^*$ by $a$, $-\tfrac a2$ and $\tfrac a2$ times the classical transformation. If the differences of the coefficients along $P_b$ and along $\Delta$ depend on the sector, a multiple of $\{S^{(1)}(\Delta),-\}$ removes $P_1$ on the families with letters of one sector, and the translation-covariant scheme removes it on every family.
\end{itemize}
\end{corollary}

\begin{proof}
(i) By Corollary \ref{cor:c11-D1local} and Proposition \ref{prop:c11-position} (i), the position part of $D_1$ is twice the sum over the pairs of the composite pair kernels along $P_b$ minus the bracket with the position part of $\hbar F_1^{(1)}(\Delta)$, by Lemma \ref{lem:c11-bracketpos}; the coefficients of $\log\varepsilon$ and the bracket with the field-counting functional commute with the translations. Theorem \ref{thm:c11-arity2} gives exactness, and the position part of each sector part of $S^{(1)}(\Delta)$ is twice the sector part of $S^{(1)}(P_b)$ (Definition \ref{def:c11-kernels}). (ii) For equal coefficients the $\theta$-terms are $2\theta$ units times $S^{(1)}(P_b)$, which is $s$-closed, and the other directions are $s$-exact after completion, with $\{sG,O\}=-s\{G,O\}$. (iii) In the covariant scheme the position part of $\hbar F_1^{(1)}(\Delta)$ is $2\hbar F_1^{(1)}(P_b)$, and $P_1$ is the sum of the composite pair kernels along $P_b$ minus the bracket with $\hbar F_1^{(1)}(P_b)$ (Corollary \ref{cor:c11-D1local}). (iv) By (ii), $P_1\cong\theta$ units times $\{S^{(1)}(P_b),-\}$ minus $\theta_P$ units times the same; Proposition \ref{prop:c9-scheme} with $\lambda_{P_b}=-2$. For sector-dependent differences $P_1$ induces the sum over the sectors of the differences times the brackets with the sector parts of $S^{(1)}(P_b)$, plus the completions, which is a multiple of $\{S^{(1)}(P_b),-\}$ on the cocycles with letters of one sector.
\end{proof}

\section{The dilatation identities in kernel form}

Corollary \ref{cor:c10-P7} assumes that on cohomology the composite term minus the bracket with the finite counterterm beyond the field-counting functional induces $-\tfrac13$ units times $\{S^{(1)}(\Delta),-\}$. By Corollary \ref{cor:c11-D1local} this is a statement about the composite pair kernels, which splits into a statement along $P_b$ and a statement about the translation-invariant part.

\begin{lemma}[One-body terms]
\label{lem:c11-onebody}
For $\theta=\theta_\lambda=\theta_V$, $\kappa_V=\kappa_\lambda=0$ and the quadratic part of $\hbar F_1^{(1)}(\Delta)$ of Definition \ref{def:c11-free}, the translation-invariant part of $-\big\{\hbar F_1^{(1)}(\Delta)-4\,\text{units}\cdot(\text{field-counting functional}),X\big\}$ on a field letter $X=\partial^\alpha\varphi$ is $(c_\varphi-2\theta\,w_X)$ units times $X$, with $c_\varphi=\tfrac23$ for $A|_V$ and $A|_{W_{\mathbb R}^\perp}$, $\tfrac13$ for $\lambda$ and $\tfrac43$ for $c$, and its position part along $b$ is $-2\theta$ units times $\partial_bX$. At $\theta=-\tfrac13$ the translation-invariant part is $\tfrac23(2+|\alpha|)$ units times $X$ for every field slot and for $c$.
\end{lemma}

\begin{proof}
$\{\int(A,A^*),X\}=X$, $\{\int(\lambda^*,\lambda),\lambda\}=-\lambda$ and $\{\int(c^*,c),c\}=-c$ (Lemma \ref{lem:c4-count}), and the bracket of the sector part with a letter of its sector is $v.X$, with translation-invariant part $2w_XX$ and position part $2\partial_bX$ (Lemma \ref{lem:c11-canonical}). At $\theta=-\tfrac13$: $\tfrac23+\tfrac23(1+|\alpha|)$, $\tfrac13+\tfrac23(\tfrac32+|\alpha|)$ and $\tfrac43+\tfrac23|\alpha|$ are all $\tfrac23(2+|\alpha|)$.
\end{proof}

\begin{definition}[The identities in kernel form]
\label{def:c11-P7}
Let $\mathfrak F$ be a family of gauge-invariant cocycles of field letters, stable under $\Lambda$, satisfying (ii) and (iii) of Definition \ref{def:c10-hyp}. The letter count is the derivation that multiplies a monomial $\int\omega\,X_1\cdots X_n$ of letters $X_a=\partial^{\alpha_a}\varphi_a$, $\varphi_a\in\{A|_V,A|_{W_{\mathbb R}^\perp},\lambda\}$, by $\sum_a(2+|\alpha_a|)$; it is defined on canonical forms, and its commutator with $\{S^{(1)}(P_b),-\}$ is $\{S^{(1)}(P_b),-\}$. The three identities on $\mathfrak F$ are:
\begin{itemize}
\item (a) the cubic part of \eqref{eq:c9-arity1} forces $\theta=\theta_\lambda=\theta_V$;
\item (b) the sum over the pairs, with the Koszul signs of \eqref{eq:c10-decomp}, of the composite pair kernels along $P_b$ induces $-\tfrac13$ units times $\{S^{(1)}(P_b),-\}$ on the cohomology of $\mathfrak F$;
\item (c) the sum over the pairs of the translation-invariant parts of the composite pair kernels along $\Delta$ induces $-\tfrac23$ units times the letter count on the cohomology of $\mathfrak F$.
\end{itemize}
\end{definition}

\begin{remark}[Dependences among the identities]
\label{rem:c11-P7dep}
(1) By Proposition \ref{prop:c11-position} (ii), pair by pair the commutator of the translation-invariant part of the composite pair kernel along $\Delta$ with $\{S^{(1)}(P_b),-\}$ is twice the composite pair kernel along $P_b$, and the commutator of the letter count with $\{S^{(1)}(P_b),-\}$ is $\{S^{(1)}(P_b),-\}$; hence (c) on $\mathfrak F$ and on $\{S^{(1)}(P_b),\mathfrak F\}$ gives (b) along $b$, and the content of (b) beyond (c) lies in the directions $b$ with $\{S^{(1)}(P_b),\mathfrak F\}$ outside $\mathfrak F$. (2) Identity (a) follows from (b) on three families with letters of one sector each, under the hypotheses of Theorem \ref{thm:c11-arity2} (Proposition \ref{prop:c11-theta-sector}); the cubic arity-one equation is then a consistency condition. (3) Identities (b) and (c) are identities between graph sums at the regulator of Definition \ref{def:c8-PT}; the numbers $-\tfrac13$ and $-\tfrac23$ are finite parts of special-edge cycles, which depend on the cutoff shape (Remark \ref{rem:c10-scheme1} (d), Proposition \ref{prop:c10-shape}). The statement independent of the scheme is the spectrum of $D_1$ at the coefficients fixed by the arity-two equation (Remark \ref{rem:c10-scheme1}).
\end{remark}

\begin{theorem}[Consequences of the identities]
\label{thm:c11-P7}
Assume (a), (b) and (c) of Definition \ref{def:c11-P7} on $\mathfrak F$.
\begin{itemize}
\item (i) $D_1$ minus four times the map of Theorem \ref{thm:c10-brst} induces $-(\tfrac13+\theta)$ units times $\{S^{(1)}(\Delta),-\}$ on the cohomology of $\mathfrak F$: the assumption of Corollary \ref{cor:c10-P7}, for the solution with sector coefficient $\theta$.
\item (ii) Under the hypotheses of Theorem \ref{thm:c11-arity2}, and if $\{S^{(1)}(P_b),O\}$ is not exact for some cocycle $O\in\mathfrak F$, $\theta=-\tfrac13$, and $D_1$ induces four times the map of Theorem \ref{thm:c10-brst}; its spectrum on $\mathfrak F$ in a fixed classical Weyl weight is twice that of the one-loop dilatation coefficients, $-8$ units times the eigenvalues of $\sum_lH_{l,l+1}$ of [B03c, (1.1)--(1.3)] on planar families (Corollary \ref{cor:c10-beisert}).
\item (iii) In the translation-covariant scheme $P_1\cong0$ on $\mathfrak F$; otherwise Corollary \ref{cor:c11-theta} (iv) applies.
\end{itemize}
Conversely, if $D_1$ minus four times the map induces $-(\tfrac13+\theta)$ units times $\{S^{(1)}(\Delta),-\}$ for equal sector coefficients, then (b) and (c) hold on $\mathfrak F$.
\end{theorem}

\begin{proof}
(i) By Corollary \ref{cor:c11-D1local}, $D_1$ minus four times the map is the sum of the composite pair kernels along $\Delta$ minus the bracket with $\hbar F_1^{(1)}(\Delta)-4$ units times the field-counting functional. Its position part along $b$ is, by (b) and Lemma \ref{lem:c11-onebody}, $2\cdot(-\tfrac13)\partial_b-2\theta\,\partial_b=-2(\tfrac13+\theta)\partial_b$ units. Its translation-invariant part is, by (c) and Lemma \ref{lem:c11-onebody}, $-\tfrac23$ units times the letter count plus $\sum_a(c_{\varphi_a}-2\theta w_{X_a})$ units, which is $-2(\tfrac13+\theta)w_O$ units, by $\sum_a(c_{\varphi_a}+\tfrac23w_{X_a})=\tfrac23\sum_a(2+|\alpha_a|)$. Together they give $-(\tfrac13+\theta)(2x\cdot\partial+2w_O)$ units, which is $-(\tfrac13+\theta)$ units times $v.$ (Lemma \ref{lem:c11-canonical}). (ii) By Theorem \ref{thm:c11-arity2} the position part of $D_1$, which is $-2(\tfrac13+\theta)$ units times $\partial_b$ on cohomology, is exact, and $\partial_bO$ is not exact for that cocycle, so $\theta=-\tfrac13$; then Theorem \ref{thm:c10-brst} and Corollary \ref{cor:c10-beisert}. (iii) Corollary \ref{cor:c11-theta} (iii). The converse reverses the two computations of (i).
\end{proof}

\begin{proposition}[The arity-two equation fixes each sector coefficient]
\label{prop:c11-theta-sector}
Assume the hypotheses of Theorem \ref{thm:c11-arity2} and of Definition \ref{def:c10-hyp}. Let $\mathfrak F_\perp$, $\mathfrak F_\lambda$, $\mathfrak F_V$ be families of gauge-invariant cocycles whose two-field parts have letters of the sector of $A|_{W_{\mathbb R}^\perp}$, of $\lambda$ and of $A|_V$ alone, each containing a cocycle $O$ for which $\{S^{(1)}(P_b),O\}$ is not exact. If (b) holds on $\mathfrak F_\perp$, then $\theta=-\tfrac13$; if it holds on $\mathfrak F_\lambda$, then $\theta_\lambda=-\tfrac13$; if it holds on $\mathfrak F_V$, then $\theta_V=-\tfrac13$. In particular (b) on the three families gives $\theta=\theta_\lambda=\theta_V=-\tfrac13$, the equality asserted in (a), and the cubic part of \eqref{eq:c9-arity1} admits the solution $\theta-\theta_V=\theta_\lambda-\theta_V=0$.
\end{proposition}

\begin{proof}
By Corollary \ref{cor:c11-D1local}, Proposition \ref{prop:c11-position} (i) and Lemma \ref{lem:c11-bracketpos}, the position part of $D_1$ is twice the sum of the composite pair kernels along $P_b$ minus the bracket with the position part of $\hbar F_1^{(1)}(\Delta)$. On the two-field part $\hbar F_1^{(1)}(\Delta)$ acts through its quadratic part; the bracket of the sector part of the sector of $X$ has the position part $2\partial_bX$ and the other sector parts have bracket $0$ with $X$, and the $\kappa$-directions are $s$-exact (Proposition \ref{prop:c11-arity1} (ii)). With (b), on the cohomology of the family of the sector with coefficient $\theta_\star\in\{\theta,\theta_\lambda,\theta_V\}$,
\[
\big(\text{position part of }D_1\text{ along }b\big)\cong-\tfrac23\,\text{units}\cdot\partial_b-2\theta_\star\,\text{units}\cdot\partial_b=-2\big(\tfrac13+\theta_\star\big)\,\text{units}\cdot\partial_b .
\]
The left side is $s$-exact on cocycles (Theorem \ref{thm:c11-arity2}) and $\partial_bO$ is not exact, so $\theta_\star=-\tfrac13$. The last statement is Proposition \ref{prop:c11-arity1} (iii); the existence of the quantisation is a hypothesis of Theorem \ref{thm:c11-arity2}.
\end{proof}

\begin{corollary}[The dilatation component without the cubic equation]
\label{cor:c11-noA}
Under the hypotheses of Proposition \ref{prop:c11-theta-sector}, if (b) holds on $\mathfrak F_\perp$, $\mathfrak F_\lambda$ and $\mathfrak F_V$, and (b) and (c) hold on $\mathfrak F$, then on the cohomology of $\mathfrak F$, $D_1$ induces four times the map of Theorem \ref{thm:c10-brst}, with spectrum $-8$ units times the eigenvalues of $\sum_lH_{l,l+1}$ on planar families, and $P_1\cong0$ in the translation-covariant scheme.
\end{corollary}

\begin{proof}
Proposition \ref{prop:c11-theta-sector} gives $\theta=\theta_\lambda=\theta_V=-\tfrac13$, and Theorem \ref{thm:c11-P7} (i)--(iii) applies.
\end{proof}

\begin{proposition}[The planar scalar sector]
\label{prop:c11-scalar}
On the part with the outputs $A|_{W_{\mathbb R}^\perp}$ of the planar scalar sector (Proposition \ref{prop:c10-scalarsector}), (b) and (c) hold as identities modulo terms whose sum over the cyclically adjacent pairs of every trace vanishes: per adjacent pair the composite pair kernel along $\Delta$ is $\big[-\tfrac16v.(X_aX_b)-\tfrac23X_aX_b\big]$ units and along $P_b$ it is $-\tfrac16$ units times $\partial_b(X_aX_b)$, so that on $O=\operatorname{tr}(\iota_{w_1}A\cdots\iota_{w_L}A)$ the composite term along $\Delta$ is $-\big(\tfrac13v.+\tfrac23L\big)O$ units, the composite term along $P_b$ is $-\tfrac13$ units times $\partial_bO$, and the sum of the translation-invariant parts is $-\tfrac43L$ units, $-\tfrac23$ units times the letter count. On the bilinears these are exact identities (Example \ref{ex:c11-bilinear}).
\end{proposition}

\begin{proof}
By Proposition \ref{prop:c10-scalarsector} (i), (e), the composite term of an adjacent pair is $-v.(X_aX_b)$ units, and the bracket with the renormalised composite two-point functional on a letter is $(\tfrac53v.X-\tfrac23X)$ units (Remark \ref{rem:c9-Pi}); by Theorem \ref{thm:c11-local} (v) their sum is the sum of the composite pair kernels, and the one-body term distributed over the two adjacent pairs of each letter gives the density per pair. For $L$ letters without derivatives the letter count is $2L$ and the translation-invariant part of $v.$ is $2L$: $-\tfrac13\cdot2L-\tfrac23L=-\tfrac43L$. Along $P_b$: Theorem \ref{thm:c9-P1}, $-2+\tfrac53=-\tfrac13$.
\end{proof}

\section{The composite kernels on the families}

Throughout this section $\mathfrak g=\mathfrak{su}(N)$ at leading order in $N$ and $O$ is a single trace, with the conventions of Section 10.7. The projection to cohomology on the two-field part is the evaluation on the free equations of motion, $|p|^2=0$ on the scalar legs, the Weyl equation on the gaugino legs and transverse on-shell gauge legs, which removes the $s$-exact terms $s(\text{antifield}\times\text{field})$; total derivatives are kept, as descendants through $\omega$. The letters are written in the components of Definition \ref{def:c8-twocomp}.

\begin{definition}[Test families]
\label{def:c11-families}
Let $n_m:=\sigma^m_{1\dot1}=(0,0,1,i)$, a null vector, and $D_n:=n_m(\partial_m+[A_m,-])$, the covariant derivative $\iota_nd_{A|_V}$ along $n$; let $Z=\iota_wA$ with $w\in W_{\mathbb R}^\perp\otimes\mathbb C$, $(w,w)=0$; let $\psi:=\lambda_1{}^i$ for a fixed $i$, the component $\alpha=1$ of $\lambda|_{S_+\otimes W}$; and let $F_{11}$ be the component $\alpha=\beta=1$ of $(\sigma^m\bar\sigma^n)_\alpha{}^\gamma\varepsilon_{\gamma\beta}F_{mn}$, $F_{mn}=\partial_mA_n-\partial_nA_m+[A_m,A_n]$, whose two-field part is a component of $dA|_V$ and whose three-field part is a component of $[A|_V\wedge A|_V]$. The test families are the spans of the following single traces and their gauge completions:
\begin{itemize}
\item the scalars: $\operatorname{tr}(\iota_{w_1}A\cdots\iota_{w_L}A)$, the bilinears, and the $\mathfrak{su}(2)$ sector of Proposition \ref{prop:c10-su2} for $L\le8$;
\item twist two: $\operatorname{tr}\big(D_n^kZ\,D_n^{S'-k}Z\big)$, $0\le k\le S'$ (Proposition \ref{prop:c10-twist});
\item the $\mathfrak{su}(1|1)$ sector: single traces of the letters $Z$ and $\psi$ of Proposition \ref{prop:c10-su11}, $L\le6$;
\item the fermion $\mathfrak{sl}(2)$ sector: $\operatorname{tr}\big(D_n^k\psi\,D_n^{S'-k}\psi\big)$;
\item two gaugino components: $\operatorname{tr}\big(D_n^k\psi^{(3)}\,D_n^{S'-k}\psi^{(4)}\big)$ with $\psi^{(3)}=\lambda_1{}^3$, $\psi^{(4)}=\lambda_1{}^4$, together with $\operatorname{tr}\big(D_n^kZ'\,D_n^{S'-k}F_{11}\big)$, $Z'$ the scalar letter with the R-symmetry weights of $\psi^{(3)}\psi^{(4)}$, with which the family closes;
\item the gluon $\mathfrak{sl}(2)$ sector: $\operatorname{tr}\big(D_n^kF_{11}\,D_n^{S'-k}F_{11}\big)$;
\item the mixed gluon--gaugino family: $\operatorname{tr}\big(D_n^kF_{11}\,D_n^{S'-k}\psi\big)$.
\end{itemize}
Twist two, the fermion and the gluon $\mathfrak{sl}(2)$ sectors are families of Proposition \ref{prop:c11-theta-sector} for the sectors of $A|_{W_{\mathbb R}^\perp}$, of $\lambda$ and of $A|_V$.
\end{definition}

The two-field part of a gauge-invariant operator with covariant derivatives receives the pairs of its two leading letters with two outputs and the pairs of a letter $A|_V$ from a covariant derivative, or from the three-field part of a field strength, with another letter, with one output (Proposition \ref{prop:c10-twist}). For the composite pair kernels these are the leg sets $A|_{W_{\mathbb R}^\perp}$ of $(A|_V,A|_{W_{\mathbb R}^\perp})$, $\lambda$ of $(A|_V,\lambda)$ and $A|_V$ of $(A|_V,A|_V)$ of Table \ref{tab:c11-fieldpairs}, all with $k=1$, with derivatives.

\begin{proposition}[Spectra from the harmonic action]
\label{prop:c11-targets}
On the two-site graded cyclic states, single traces of two letters with $S'$ derivatives $D_n$ in total, the eigenvalues of $\sum_lH_{l,l+1}$ of [B03c, (1.1)--(1.3)], computed from the harmonic action [B03c, (4.27)--(4.29)] on the oscillator letters [B03c, (4.9)], are:
\begin{itemize}
\item (i) scalars $D_n^kZ$: $4h(S)$, $S\le S'$ even [B03c, (2.31)];
\item (ii) gluons $D_n^kF_{11}$: $4h(S+2)$, $S\le S'$ even, so that $\operatorname{tr}(F_{11}F_{11})$ has the Konishi value $6$;
\item (iii) one gaugino component $D_n^k\psi$: $4h(S+1)$, $S\le S'$ odd;
\item (iv) two gaugino components $\psi^{(3)}$, $\psi^{(4)}$ with $Z'F_{11}$: $0,6$ at $S'=0$; $0,6,6,6$ at $S'=1$; $0,6,6,6,6,\tfrac{25}3$ at $S'=2$; $0$, four times $6$ and three times $\tfrac{25}3$ at $S'=3$; $0$, four times $6$, four times $\tfrac{25}3$ and $\tfrac{49}5$ at $S'=4$;
\item (v) $(F_{11},\psi)$: $6$; $6,6$; $6,6,\tfrac{25}3$; $6,6,\tfrac{25}3,\tfrac{25}3$; $6,6,\tfrac{25}3,\tfrac{25}3,\tfrac{49}5$ at $S'=0,\dots,4$.
\end{itemize}
On the $\mathfrak{su}(1|1)$ sector the two-site Hamiltonian is one minus the graded exchange (Proposition \ref{prop:c10-su11}), with the cyclic spectra of Table \ref{tab:c11-su11}.
\end{proposition}

\begin{proof}
Exact evaluation of [B03c, (4.27)--(4.29)], with the harmonic numbers of the coefficients $c_{n,0,0}=h(\tfrac n2)$, on bases of the families closed under $H_{12}$, with the graded exchange for $H_{21}$ and the projection to graded cyclic states, and exact diagonalisation. The same evaluation gives [B03c, (2.31)] on twist two and the values $0$, twenty times, and $6$ on the scalar bilinears, which fixes the central charge per site, the signs of the fermionic oscillators and the graded cyclicity.
\end{proof}

\begin{tabular}{l|l}
$(L,M)$ & eigenvalues on the graded cyclic states of $\operatorname{tr}(Z^{L-M}\psi^M)$\\
\hline
$M\le1$ & $0$\\
$(3,2)$, $(3,3)$ & $6$\\
$(4,2)$, $(4,3)$ & $4$\\
$(5,2)$, $(5,3)$ & $5\pm\sqrt5$\\
$(5,4)$, $(5,5)$ & $10$\\
$(6,2)$ & $2$, $6$\\
$(6,3)$ & $2$, $6$, $8$, $8$\\
$(6,4)$ & $8$, $8$, $8$\\
$(6,5)$ & $8$\\
\end{tabular}
\tabcaption{The eigenvalues of $\sum_lH_{l,l+1}$ on the $\mathfrak{su}(1|1)$ sector, $L\le6$; the sectors $(2,2)$, $(4,4)$, $(6,6)$ have no graded cyclic state.\label{tab:c11-su11}}

\begin{theorem}[The composite pair kernels on the families]
\label{thm:c11-families}
\begin{itemize}
\item (i) (The two-site density.) For every ordered adjacent pair $(X_a,X_b)$ of letters among $A|_{W_{\mathbb R}^\perp}$, $\lambda|_{S_+\otimes W}$, $\lambda|_{S_-\otimes W^*}$ without derivatives, on the free shell, and for two letters $A|_{W_{\mathbb R}^\perp}$ also off the shell and for all their components, the planar composite pair kernels have, on the same word $X_aX_b$, the translation-invariant part $-\tfrac43$ units times $X_aX_b$ along $\Delta$ and the value $-\tfrac16$ units times $\partial_b(X_aX_b)$ along $P_b$; on the exchanged word and on the trace they vanish.
\item (ii) On the two-field parts of the families of Definition \ref{def:c11-families}, on the free shell, the sum over the pairs of the translation-invariant parts of the composite pair kernels along $\Delta$ is the multiple of the identity of Table \ref{tab:c11-results}, which is $-\tfrac23$ units times the letter count, and the sum of the composite pair kernels along $n$ is $-\tfrac13$ units times $\partial_n$; on twist two for $S'\le4$ the sum along every $b\in V$ is $-\tfrac13$ units times $\partial_b$.
\item (iii) Hence (b) along $n$ and (c) of Definition \ref{def:c11-P7} hold on each family in the ranges of Table \ref{tab:c11-results}, and (b) along every $b$ holds on twist two for $S'\le4$.
\end{itemize}
\end{theorem}

\begin{proof}
Exact evaluation of the special-edge composite cycles by Theorem \ref{thm:c11-simplex}, Propositions \ref{prop:c11-genfun} and \ref{prop:c11-k1}, with the colour words reduced to the planar adjacent form of Lemma \ref{lem:c10-colour} (iv), on the leg sets of Table \ref{tab:c11-fieldpairs} with the derivatives of the families, the outputs evaluated on the free shell (Remark \ref{rem:c11-method}). On the families with derivatives the outputs depend on the output momenta through $(n,p)$, every derivative being contracted with $n$ and $(n,n)=0$; the sums are diagonal in the bases of the families. (iii) is (i), (ii) and the values of the letter count: $2L$ on $L$ letters without derivatives, $4+S'$ on twist two and on two gaugino letters, $5+S'$ on $Z'F_{11}$ and on $(F_{11},\psi)$, $6+S'$ on two gluon letters, a letter $F_{11}$ counting $3$ through its two-field part $dA|_V$.
\end{proof}

\begin{tabular}{l|l|l|l}
family & range & sum of the translation-invariant parts along $\Delta$ & spectrum of $D_1$ at $\theta=\theta_\lambda=\theta_V=-\tfrac13$, in units\\
\hline
scalars & bilinears; $\mathfrak{su}(2)$, $L\le8$ & $-\tfrac43L$ & $0$, $-48$; $-8$ times the Heisenberg spectra\\
twist two & $S'\le6$ & $-\tfrac23(4+S')$ & $0$, $-48$, $-\tfrac{200}3$, $-\tfrac{392}5$: $-32h(S)$, $S\le S'$ even\\
$\mathfrak{su}(1|1)$ & $L\le6$ & $-\tfrac43L$ & $-8$ times the $\mathfrak{su}(1|1)$ spectra\\
fermion $\mathfrak{sl}(2)$ & $S'\le5$ & $-\tfrac23(4+S')$ & $-48$, $-\tfrac{200}3$, $-\tfrac{392}5$: $-32h(S+1)$, $S\le S'$ odd\\
$\psi^{(3)}\psi^{(4)}$ with $Z'F_{11}$ & $S'\le2$ & $-\tfrac23(4+S')$ on $\psi^{(3)}\psi^{(4)}$, $-\tfrac23(5+S')$ on $Z'F_{11}$ & $0,-48$; $0,-48,-48,-48$; $0$, four times $-48$, $-\tfrac{200}3$\\
gluon $\mathfrak{sl}(2)$ & $S'\le2$ & $-\tfrac23(6+S')$ & $-48$, $-\tfrac{200}3$: $-32h(S+2)$, $S\le S'$ even\\
$(F_{11},\psi)$ & $S'\le4$ & $-\tfrac23(5+S')$ & at $S'=4$: $-48$, $-48$, $-\tfrac{200}3$, $-\tfrac{200}3$, $-\tfrac{392}5$\\
\end{tabular}
\tabcaption{The composite pair kernels and the spectrum of $D_1$ on the families: the third column is identity (c), the fourth is $-8$ units times the eigenvalues of Proposition \ref{prop:c11-targets}, of Proposition \ref{prop:c10-su2} on the $\mathfrak{su}(2)$ sector and of Table \ref{tab:c11-su11} on the $\mathfrak{su}(1|1)$ sector.\label{tab:c11-results}}

\begin{corollary}[$D_1$ and $P_1$ on the families]
\label{cor:c11-D1fam}
Under the hypotheses of Definition \ref{def:c10-hyp}, on the cohomology of the two-field part of every family of Definition \ref{def:c11-families}, with every free direction of Definition \ref{def:c11-free},
\[
D_1\cong4\,(\text{the map})-\sum_a\big(\theta_{X_a}+\tfrac13\big)\,\text{units}\cdot v.X_a,\qquad P_1\cong-\sum_a\big(\theta_{P,X_a}+\tfrac13\big)\,\text{units}\cdot\partial_bX_a,
\]
with the map of Theorem \ref{thm:c10-brst}, where $\theta_{X_a}$ is $\theta$, $\theta_\lambda$ or $\theta_V$ and $\theta_{P,X_a}$ is $\theta_P$, $\theta_{P,\lambda}$ or $\theta_{P,V}$ according to the sector of the letter $X_a$, $v.X_a$ and $\partial_bX_a$ act on the letter $X_a$ alone, and the directions $Q\int(\iota_xA|_V,c^*)$, $Q\int(\rho(x)\lambda^*,\lambda^*)$ contribute $s$-exact terms. Under the hypotheses of Theorem \ref{thm:c11-arity2}, twist two gives $\theta=-\tfrac13$, the fermion $\mathfrak{sl}(2)$ sector $\theta_\lambda=-\tfrac13$ and the gluon $\mathfrak{sl}(2)$ sector $\theta_V=-\tfrac13$ (Proposition \ref{prop:c11-theta-sector}); then $D_1$ induces four times the map of Theorem \ref{thm:c10-brst} on every family, with the spectra of Table \ref{tab:c11-results}, $-8$ units times the eigenvalues of $\sum_lH_{l,l+1}$, which is twice the one-loop anomalous dimensions; and $P_1\cong0$ in the translation-covariant scheme, while for equal coefficients $\theta_P$ along $P_b$, $P_1\cong(\theta-\theta_P)$ units times $\partial_b$ with $\theta=-\tfrac13$.
\end{corollary}

\begin{proof}
Corollary \ref{cor:c11-D1local} and Theorem \ref{thm:c11-families}: on the two-field part $\hbar F_1^{(1)}(\Delta)$ acts through its quadratic part, the bracket of the sector part of the sector of $X_a$ with $X_a$ being $v.X_a$ and the others vanishing. As in the proof of Theorem \ref{thm:c11-P7} (i), letter by letter, the position part of the composite pair kernels and of the bracket is $-2(\tfrac13+\theta_{X_a})\partial_b$ units and the translation-invariant part is $-2(\tfrac13+\theta_{X_a})w_{X_a}$ units; along $P_b$ the composite pair kernels give $-\tfrac13$ units times $\partial_b$ and the bracket with $\hbar F_1^{(1)}(P_b)$ gives $-\theta_{P,X_a}$ units times $\partial_bX_a$. The $\kappa$-directions are $s$-exact (Proposition \ref{prop:c11-arity1} (ii)). The rest is Proposition \ref{prop:c11-theta-sector}, Theorem \ref{thm:c11-P7} and Corollary \ref{cor:c11-theta} (iii), (iv).
\end{proof}

\begin{remark}[Method of computation]
\label{rem:c11-method}
Each special-edge composite cycle is evaluated exactly: the Gaussian loop integral and the integrals of \eqref{eq:c11-simplex} over the ordered subsimplices of the region where $\alpha_\sigma$ is the smallest coordinate, the letter derivatives through Proposition \ref{prop:c11-genfun}, the position through the injected momentum $r$ (Proposition \ref{prop:c11-position}), the chains with $k=1$ through Proposition \ref{prop:c11-k1}, the fibre algebra in explicit realisations of $V_{10}$ and $\Sigma$ over $\mathbb Q(i)$, and the colour words by Lemma \ref{lem:c10-colour}, with its planar adjacent rule also checked in realisations of $\mathfrak{su}(2)$, $\mathfrak{su}(3)$ and $\mathfrak{gl}(3)$. In this evaluation the factor $(-1)^{[c<\sigma]}$ of Definition \ref{def:c11-sp} is part of the chain convention when the special edge ends at $X_b$ and is applied as the factor $-1$ when the special edge joins two vertices of $I$. A second evaluation works in momentum space with the kernel form of Lemma \ref{lem:c11-kernels} (ii), $v.$ acting on the Fourier components as an Euler operator plus constants, Feynman parameters for the edges $P$ and the special Schwinger time integrated last, and applies the sign of Definition \ref{def:c11-sp} as written. The two evaluations give the same off-shell entries with at most two outputs and no letter derivatives, for the pairs of bosonic letters and for the pairs with a gaugino letter or output, whose values lie in $\mathbb Q+\mathbb Q\log2+\mathbb Q\log3$ for $k\le2$, and the same family matrices on the scalars, on twist two for $S'\le4$, on the $\mathfrak{su}(1|1)$ sector, on the fermion $\mathfrak{sl}(2)$ sector for $S'\le5$ and on the gluon $\mathfrak{sl}(2)$ sector for $S'\le2$; the first evaluation extends twist two to $S'\le6$, and the second evaluation gives the two families with $F_{11}$ and a gaugino letter. A third evaluation, a Taylor expansion in momentum space at generic null leg momenta, gives the sums on twist two for $S'\le4$ along a generic $b$. Every integral is a four-dimensional Schwinger integral with the cutoff $\varepsilon$ on every edge.
\end{remark}

\section{The two-body kernels of $Q_1$, $S_1$ and $K_1$}

For $x$ in a homogeneous basis of $\mathfrak{psl}(4|4)$, $\lambda_x$ is its $\mathrm{ad}_\Delta$ eigenvalue, $-1$ on $\Sigma$, $+1$ on $\Sigma^*$ and $+2$ for $K_b$, and the degree of $x$ in the position is $0$, $1$ and $2$ (Section 10.9). The two-body kernel of a pair of letters along $x$ is that of Proposition \ref{prop:c10-arity1}: the limit as $\varepsilon\to0$ of the local parts of the cycles through the pair with exactly one defect insertion, a special edge, the vertex $S^{(1)}(x)$ on the cycle, or the arity-one counterterm vertex of the operator. This section gives the lines that $S^{(1)}(x)$ forms on a cycle along $\Sigma$, $\Sigma^*$ and $K_b$, the output bounds, the relations of $S_1$ to $Q_1$ and of $K_1$ to $P_1$, the renormalisation freedom in the scheme of Sections 11.3--11.5, and the comparison of $Q_1$ with the one-loop differential of the holomorphic twist.

\begin{lemma}[The lines along $\Sigma$, $\Sigma^*$ and $K_b$]
\label{lem:c11-lines}
\begin{itemize}
\item (i) For $\xi\in\Sigma$ the quadratic part of $S^{(1)}(\xi)$ is $\int\mathrm{dvol}\big(-(\Gamma(\xi,\lambda),A^*)+\tfrac12(\rho(dA)\xi,\lambda^*)\big)$ (Proposition \ref{prop:c7-susy}). With both legs on edges $P$, its term $-\int(\Gamma(\xi,\lambda)|_V,A^*|_V)$ has the leg $\lambda$ on the component $\lambda\otimes\lambda$ of $P$ and the leg $A^*|_V$ on the component $A^*|_V\otimes c$: an odd line joining a slot $\lambda$ to a slot $c$, linear in $\xi$ and free of the position. The terms $-\int(\Gamma(\xi,\lambda)|_{W_{\mathbb R}^\perp},A^*|_{W_{\mathbb R}^\perp})$ and $\tfrac12\int(\rho(dA)\xi,\lambda^*)$, with a leg $A^*|_{W_{\mathbb R}^\perp}$ or $\lambda^*$, enter a cycle at the special edge.
\item (ii) For $\eta\in\Sigma^*$, $S^{(1)}(\eta)$ is $S^{(1)}(\xi)$ with $\xi$ replaced by the value $\rho(x)\eta$ of the conformal Killing spinor, plus the translation-invariant term $-\int(\rho(A|_{W_{\mathbb R}^\perp})\eta,\lambda^*)$ \eqref{eq:c2-S1S}. Its line on a cycle is the line of (i) along $\rho(x)\eta$, linear in the position; the last term, with its leg $\lambda^*$, enters a cycle at the special edge.
\item (iii) For $K_b$, with the vector field $v(x)=(x,x)b-2(b,x)x$, $S^{(1)}(K_b)$ forms on a cycle the composite line from $c$ to $A|_V$ with the kernel $P\star\,v.\,P\star=v.(P\star)^2+[P\star,v.]P\star$, quadratic in the position, at a finite endpoint $T$ (Definition \ref{def:c8-K} and the remark after Theorem \ref{thm:c8-K}).
\end{itemize}
In each case every letter, $A^*|_{W_{\mathbb R}^\perp}$, $\lambda^*$ and $c^*$ included, is active (Lemma \ref{lem:c10-active1}).
\end{lemma}

\begin{proof}
The components of $P$ are those of Proposition \ref{prop:c3-P}; $\Gamma(\xi,\lambda)\in V_{10}$ pairs with $A^*$ through its parts in $V$ and in $W_{\mathbb R}^\perp$. (ii) is \eqref{eq:c2-S1S}, the first two terms being those of \eqref{eq:c2-S1Q} at the spinor $\rho(x)\eta$. (iii) Definition \ref{def:c7-vdot} for the vector field of $K_b$ and Lemma \ref{lem:c11-kernels} (iii), (iv) with this $v$.
\end{proof}

\begin{proposition}[Output bounds along $\Sigma$, $\Sigma^*$ and $K_b$]
\label{prop:c11-oddbound}
Under (iii) of Definition \ref{def:c10-hyp} along $x$, the two-body kernel of $(X_a,X_b)$ is a local functional of the pair degree with:
\begin{itemize}
\item (i) along $\xi\in\Sigma$: Weyl weight $w_{X_a}+w_{X_b}+\tfrac12$, no position, parity $|X_a|+|X_b|+1$, and at most $w_{X_a}+w_{X_b}+\tfrac12$ plus the pair degree output fields, as for the second-order operator on pairs of letters of [BK, (3.7)];
\item (ii) along $\eta\in\Sigma^*$: Weyl weight $w_{X_a}+w_{X_b}-\tfrac12$ with the position at weight $-1$, at most linear in the position, parity $|X_a|+|X_b|+1$, and at most $w_{X_a}+w_{X_b}+\tfrac12$ plus the pair degree output fields;
\item (iii) along $K_b$: Weyl weight $w_{X_a}+w_{X_b}-1$ with the position at weight $-1$, at most quadratic in the position, parity $|X_a|+|X_b|$, and at most $w_{X_a}+w_{X_b}+1$ plus the pair degree output fields.
\end{itemize}
\end{proposition}

\begin{proof}
Proposition \ref{prop:c10-bound1} with $\lambda_x=-1,+1,+2$ and the degrees $0,1,2$ in the position.
\end{proof}

\begin{proposition}[The special supercharges and the supercharges]
\label{prop:c11-S}
Let $\eta\in\Sigma^*$ and $b\in V$, so that $\rho(b)\eta\in\Sigma$.
\begin{itemize}
\item (i) Graph by graph, the position part along $b$ of the two-body kernel of $(X_a,X_b)$ along $\eta$ is the two-body kernel of $(X_a,X_b)$ along $\rho(b)\eta$; its translation-invariant part is computed with $S^{(1)}(\eta)$ at the spinor $\rho(x-x_O)\eta$, $x_O$ the point of the density, together with the term $-\int(\rho(A|_{W_{\mathbb R}^\perp})\eta,\lambda^*)$.
\item (ii) Under the hypotheses of Proposition \ref{prop:c7-select} and Corollary \ref{cor:c8-homotopy}, on cohomology, the position part along $b$ of $S_1$ along $\eta$ is $Q_1$ along $\rho(b)\eta$ minus the commutator of $\{S^{(1)}(\eta),-\}$ with $P_1$.
\item (iii) On a family on which $P_1\cong0$, in particular on the families of Definition \ref{def:c11-families} in the translation-covariant scheme at $\theta=\theta_\lambda=\theta_V=-\tfrac13$, the position part of $S_1$ along $\eta$ is $Q_1$ along $\rho(b)\eta$ on cohomology, and $S_1$ along $\eta$ is fixed by $Q_1$ along $\Sigma$ up to the term of Proposition \ref{prop:c7-select} linear in ${\not d}(\rho(x)\eta)=4\eta$. On the scalar bilinears this term is absent (Corollary \ref{cor:c7-S}), and $S_1$ along $\eta$ is, on cohomology, the combination of the functionals of Corollary \ref{cor:c7-S} with the coefficients of $Q_1$ along $\Sigma$ in Corollary \ref{cor:c7-Q}.
\end{itemize}
\end{proposition}

\begin{proof}
(i) By Lemma \ref{lem:c11-lines} (ii), $S^{(1)}(\eta)=\int x^\mu s_\mu+\int s_0$, where $\int s_\mu$ is $S^{(1)}$ along the constant spinor $\rho(e_\mu)\eta\in\Sigma$ and $s_0=-(\rho(A|_{W_{\mathbb R}^\perp})\eta,\lambda^*)$. Insert $x^\mu=x_O^\mu+(x-x_O)^\mu$ and argue as in the proof of Proposition \ref{prop:c11-position} (i): on plane waves the factor $x^\mu$ is $-i\partial/\partial r^\mu$ of the momentum injected at the vertex, and its action on the momentum delta function is $-i$ times the graph with $\int s_\mu$; the same holds for the special edges, the vertices $S^{(1)}(\eta)$ on the cycle and the counterterm vertices, which are the defect insertions of Proposition \ref{prop:c10-arity1}. (ii) Corollary \ref{cor:c8-homotopy} (ii) at $(\eta,P_b)$: by Proposition \ref{prop:c2-special}, $[P_b,\eta]=\rho(b)\eta$, so $[\eta,P_b]=-\rho(b)\eta$, and $S^{(2)}(\eta,P_b)=0$ (Definition \ref{def:c2-S2}). At order $\hbar$ the identity states that the commutator of $S_1$ along $\eta$ with $\{S^{(1)}(P_b),-\}$, plus the commutator of $\{S^{(1)}(\eta),-\}$ with $P_1$, minus $Q_1$ along $\rho(b)\eta$, is a graded commutator with $\Lambda^{(0)}$, which induces $0$ on cohomology; at order $\hbar^0$ the same identity is the statement of (i) for the classical transformation, which fixes the signs. By Proposition \ref{prop:c10-bound1}, $S_1$ along $\eta$ is at most linear in the position, and Lemma \ref{lem:c11-canonical} identifies the first commutator with its position part. (iii) By Proposition \ref{prop:c7-select} the representative of $S_1$ along $\eta$ at a point $x_0$ is a term linear in the value $\rho(x_0)\eta$ plus a term linear in $4\eta$; the first has the position part given by its value at $\rho(b)\eta$, and as $\eta$ and $b$ vary, $\rho(b)\eta$ spans $\Sigma$, so by (ii) it is the representative of $Q_1$ evaluated at $\rho(x_0)\eta$ on cohomology.
\end{proof}

\begin{proposition}[Renormalisation freedom of the odd components and of $K_1$]
\label{prop:c11-freedom}
\begin{itemize}
\item (i) A field redefinition, a closed direction and the outer direction $G_B$ (Proposition \ref{prop:c8-laws} (a), (c)) change each component by minus the bracket with a local functional, a derivation of the letters: they change the one-body kernels and leave the two-body kernels unchanged. $G_B$ changes $Q_1$ along $S_+\otimes W$ by $-a$ and along $S_-\otimes W^*$ by $+a$ times the classical transformation, $S_1$ along $S_+\otimes W^*$ by $+a$ and along $S_-\otimes W$ by $-a$ times it, and leaves $K_1$ unchanged.
\item (ii) The finite counterterms $\hbar F_1^{(0)}$ and $\hbar F_1^{(1)}(x)$ enter the one-body kernels (Remark \ref{rem:c10-scheme1} (a), (b)); the sector coefficients of Definition \ref{def:c11-free} enter the components along $\Delta$ and $P_b$.
\item (iii) A finite part $\psi'$ of the operators changes the component along $x$ by $[\{S^{(1)}(x),-\},\psi']$ (Proposition \ref{prop:c8-laws} (b)). For $\psi'=\tfrac a2\{S^{(1)}(\Delta),-\}$ this is $\tfrac a2\lambda_x\{S^{(1)}(x),-\}$ (Proposition \ref{prop:c9-scheme}): $+a$, $-\tfrac a2$, $+\tfrac a2$ times the classical transformation along $K_b$, $\Sigma$, $\Sigma^*$, a change of the one-body kernels. A finite part $\psi'$ that vanishes on $O$ changes the component of $O$ along $x$ by $-\psi'(\{S^{(1)}(x),O\})$; in this way the finite parts of the operators of the Weyl weight of $\{S^{(1)}(x),O\}$ enter the two-body kernels.
\item (iv) The cutoff shape enters the finite parts of the special-edge cycles (Proposition \ref{prop:c10-shape}, Lemma \ref{lem:c10-special}), the composite pair kernels and the finite vertex \eqref{eq:c10-finitesp}.
\item (v) In the translation-covariant scheme at $\theta=\theta_\lambda=\theta_V=-\tfrac13$, $P_1\cong0$ on the families of Definition \ref{def:c11-families} (Corollary \ref{cor:c11-D1fam}), and the condition $P_1=0$ of Definition \ref{def:c8-radial} holds there with $\psi'=0$: the components along $K_b$, $\Sigma$ and $\Sigma^*$ are those of the scheme. For equal coefficients $\theta_P\ne\theta$ along $P_b$, the finite part with $a=(\theta-\theta_P)$ units gives $P_1\cong0$ and changes the components along $K_b$, $\Sigma$, $\Sigma^*$ by $a$, $-\tfrac a2$, $\tfrac a2$ times the classical transformation (Corollary \ref{cor:c11-theta} (iv)).
\item (vi) The class $[Q^\xi_1]$ on the cohomology of $Q^\xi_0$, its kernel, image and rank are invariant under (i) and (iii) (Theorem \ref{thm:c8-twist}).
\end{itemize}
\end{proposition}

\begin{proof}
(i) Proposition \ref{prop:c8-laws}; the bracket with a local functional is a derivation in its second argument, so it acts letter by letter. (ii) Remark \ref{rem:c10-scheme1} and Definition \ref{def:c8-components}. (iii) Proposition \ref{prop:c8-laws} (b) and Proposition \ref{prop:c9-scheme}, with $[\{S^{(1)}(x),-\},\psi'](O)=\{S^{(1)}(x),\psi'(O)\}-\psi'(\{S^{(1)}(x),O\})$. (iv) Remark \ref{rem:c10-scheme1} (d). (v) Corollaries \ref{cor:c11-D1fam} and \ref{cor:c11-theta} (iv). (vi) Theorem \ref{thm:c8-twist}.
\end{proof}

\begin{corollary}[$K_1$ in the translation-covariant scheme]
\label{cor:c11-K}
Under the hypotheses of Theorem \ref{thm:c8-adjoint}, on a family of Definition \ref{def:c11-families} at the origin of one length, in the translation-covariant scheme at $\theta=\theta_\lambda=\theta_V=-\tfrac13$, $K_1$ along $K_{\bar b}$ equals $\hbar\big[\{S^{(1)}(K_{\bar b}),-\},\mathrm{Gram}_{\rm tree}^{-1}\mathrm{Gram}_{\rm 1}\big]$, the statement $K=P^\dagger$ of Corollary \ref{cor:c8-gram}. For equal coefficients $\theta_P$ along $P_b$, $K_1$ along $K_{\bar b}$ is $(\theta-\theta_P)$ units times $\{S^{(1)}(K_{\bar b}),-\}$ plus the same commutator, with the radial Gram matrix of Definition \ref{def:c8-inner} in that scheme.
\end{corollary}

\begin{proof}
Corollary \ref{cor:c11-D1fam} gives $P_1\cong0$, resp. $P_1\cong(\theta-\theta_P)$ units times $\{S^{(1)}(P_b),-\}$ with a real coefficient; Corollary \ref{cor:c8-gram} gives the two relations.
\end{proof}

\begin{proposition}[The chiral-primary bilinears along a rank-one supercharge]
\label{prop:c11-plane}
Let $\xi\in S_+\otimes W$ be of rank one, $\xi_\alpha{}^i=u_\alpha z^i$, and let $w_1,w_2\in W_{\mathbb R}^\perp\otimes\mathbb C$ satisfy $\rho(w_1)\xi=\rho(w_2)\xi=0$. Then $(w_1,w_2)=0$, and under the hypotheses of Proposition \ref{prop:c7-select} the component along $\xi$ of $O=\int\omega\,(\iota_{w_1}A,\iota_{w_2}A)$ vanishes at every order at which they hold: its classical part $\{S^{(1)}(\xi),O\}$ and $Q_1(O)$ are $0$, in every scheme in which the counterterms are equivariant for translations, $\mathfrak{so}(V)$, $\mathfrak{so}(W_{\mathbb R}^\perp)$ and $\mathfrak g$.
\end{proposition}

\begin{proof}
$\rho(w_1)\rho(w_2)+\rho(w_2)\rho(w_1)=2(w_1,w_2)$ applied to $\xi\ne0$ gives $(w_1,w_2)=0$. By Corollary \ref{cor:c7-Q}, for $(w_1,w_2)=0$ the component is a multiple of $\int\omega\,\big[(\iota_{w_1}A,(\rho(w_2)\xi,\lambda))+(\iota_{w_2}A,(\rho(w_1)\xi,\lambda))\big]$, which vanishes; the classical component is the same functional (proof of Corollary \ref{cor:c7-Q}).
\end{proof}

\begin{remark}[The pair entries of the holomorphic twist]
\label{rem:c11-twist}
The twisted one-loop differential is $\hbar Q^\xi_1(O)=\hbar L_1(O)^{(0)}+Q_1(O)-\tfrac12\hbar L_1(O)^{(2)}(\xi,\xi)-\cdots$ (Definition \ref{def:c8-twisted}), so its pair rule [BK, (3.7)--(3.8)] corresponds to the two-body kernels of $Q_1$ along $\xi$ and of the arity-two component along $(\xi,\xi)$. By Proposition \ref{prop:c7-select}, applied once for each argument, a component of arity $k$ along $\xi$ of an operator of Weyl weight $w$ and degree $d$ has representatives of Weyl weight $w+\tfrac k2$ and degree $d+1-k$. [BK, (3.51)] relates the holomorphic letters to the fields up to field redefinitions; with $\xi_\alpha{}^i=u_\alpha z^i$, $u'\in S_+$ with $u'^\alpha u_\alpha=1$, and the representatives $\gamma\leftrightarrow A_{ij}z^j$ (Weyl weight $1$, degree $0$), $\beta\leftrightarrow u'^\alpha\lambda_\alpha{}^i$ modulo $z^i$ $(\tfrac32,0)$, $\partial_{\dot\alpha}c\leftrightarrow\bar\lambda_{\dot\alpha i}z^i$ $(\tfrac32,0)$, $c\leftrightarrow c$ $(0,1)$, $b\leftrightarrow u'^\alpha F_{mn}(\sigma^m\bar\sigma^n)_\alpha{}^\beta u'_\beta$ $(2,0)$ or a component of $\lambda^*$ $(\tfrac52,-1)$, and $\partial_{\dot\alpha}\leftrightarrow u'^\alpha\sigma^m_{\alpha\dot\alpha}D_m$ $(1,0)$, the pair entries of [BK, (3.64)--(3.69)] lie in the following arities:

\begin{tabular}{l|c|c|c}
entry & Weyl weight and degree of the pair & of the output & arity\\
\hline
$b\,c$ [BK, (3.64)] & $(2,1)$ & $(3,0)$ & $2$\\
$b\,b$ [BK, (3.65)] & $(4,0)$ & $(\tfrac92,0)$ & $1$\\
$\beta\,\gamma$ [BK, (3.66)] & $(\tfrac52,0)$ & $(3,0)$ & $1$\\
$\beta\,\beta$ [BK, (3.67)] & $(3,0)$ & $(\tfrac72,0)$ & $1$\\
$b\,\gamma$ [BK, (3.68)] & $(3,0)$ & $(\tfrac72,0)$ & $1$\\
$\beta\,b$ [BK, (3.69)] & $(\tfrac72,0)$ & $(4,0)$ & $1$\\
\end{tabular}
\tabcaption{The arities of the one-loop pair entries of the holomorphic twist, with $b$ represented by the self-dual field strength; with $b$ represented by a component of $\lambda^*$ the entry $b\,c$ has arity $1$.\label{tab:c11-twist}}

The entry $\beta\gamma$, the Konishi entry, is the two-body kernel of $Q_1$ along $\xi$ on the pair $(\lambda|_{S_+\otimes W},A|_{W_{\mathbb R}^\perp})$ with two outputs $\lambda|_{S_-\otimes W^*}$: the coefficient of $\int\omega\,(\bar\lambda_{\dot\alpha i}z^i,\bar\lambda^{\dot\alpha}{}_jz^j)$ in $Q_1$ of $\int\omega\,(u'^\alpha\lambda_\alpha{}^i,A_{ij}z^j)$. The classical component of this pair along $\xi$ is a number $\kappa$ times a cubic functional of the scalars, and a finite part of the operators of Weyl weight $3$ that adds $y\int\omega\,(\bar\lambda_{\dot\alpha i}z^i,\bar\lambda^{\dot\alpha}{}_jz^j)$ to that cubic functional changes the entry by $-\kappa y$ (Proposition \ref{prop:c11-freedom} (iii)), the two functionals having the same Weyl weight and quantum numbers. The pair entries are therefore compared through the class $[Q^\xi_1]$ on the cohomology of $Q^\xi_0$ (Theorem \ref{thm:c8-twist}), transported to the holomorphic letters through [ESW, Thm 10.9] (remark after Theorem \ref{thm:c8-twist}). Proposition \ref{prop:c11-plane} is the exact untwisted form, for two letters, of [BK, (3.77)], $Q_1(\gamma^I\gamma^J)=0$: for $\rho(w_a)\xi=0$, $\iota_{w_a}A$ is a combination of the components $A_{ij}z^j$, the letters $\gamma$.
\end{remark}

\begin{remark}[The data that complete the one-loop components]
\label{rem:c11-next}
The two-body kernels of $Q_1$ along $\xi\in\Sigma$ are the cycles of Proposition \ref{prop:c10-arity1} with the odd line of Lemma \ref{lem:c11-lines} (i) or with the special edge and the vertex $S^{(1)}(\xi)$ on a leg; their evaluation by the simplex formulas of Theorem \ref{thm:c11-simplex}, in the scheme of Sections 11.3--11.5 and with the signs of Definition \ref{def:c11-sp}, is the next computation: first on the Konishi pair of Remark \ref{rem:c11-twist}, on the pairs of Table \ref{tab:c11-twist} and on the scalar bilinears of Corollary \ref{cor:c7-Q}. With them, $S_1$ follows by Proposition \ref{prop:c11-S} up to the term linear in $4\eta$, and $K_1$ by Corollary \ref{cor:c11-K} with the radial Gram matrix of the two-point graphs of Definition \ref{def:c8-2pt}. The cubic part of \eqref{eq:c9-arity1}, with the cubic part of the composite part from the composite loops of $I$ with three legs, is the consistency condition of Proposition \ref{prop:c11-theta-sector}, and the one-loop quantisation through arity two along $(\Delta,P_b)$ is the hypothesis of Theorem \ref{thm:c11-arity2}. The pole identity and the logarithmic covariance of Theorem \ref{thm:c11-local} (iii) are identities for the composite cycles through the pairs and the arity-one counterterms, tabulated with the derivatives of Table \ref{tab:c11-counts}; the gauge completion of the families couples the sector coefficients through the three-field parts.
\end{remark}

\begin{remark}[Two-component forms]
\label{rem:c11-twocomp}
With the components of Definition \ref{def:c8-twocomp}:
\begin{itemize}
\item Letters. $A|_V$, $A|_{W_{\mathbb R}^\perp}$, $\lambda|_{S_+\otimes W}$, $\lambda|_{S_-\otimes W^*}$ are $A_m$, $A_{ij}$ (the $\mathbf 6$ of $SU(4)_R$), $\lambda_\alpha{}^i$ (undotted, $\mathbf 4$), $\bar\lambda_{\dot\alpha i}$ (dotted, $\bar{\mathbf 4}$), with the antifields $A^*_m$, $A^*_{ij}$, $\lambda^{*\alpha}{}_i$, $\bar\lambda^{*\dot\alpha i}$, $c^*$ and the degrees $1,0,0,0,0,-1,-1,-1,-1,-2$ for $c,A_m,A_{ij},\lambda_\alpha{}^i,\bar\lambda_{\dot\alpha i},A^*_m,A^*_{ij},\lambda^{*\alpha}{}_i,\bar\lambda^{*\dot\alpha i},c^*$; $(\lambda,\lambda^*)=(\lambda_\alpha{}^i,\lambda^{*\alpha}{}_i)-(\bar\lambda_{\dot\alpha i},\bar\lambda^{*\dot\alpha i})$ and $(u,w)=\tfrac14u^{ij}w_{ij}$ on $W_{\mathbb R}^\perp$.
\item Composite kernels. The internal vertex along $\Delta$ is $\int\big((2E+2)A_m,A^*_m\big)=2\int x_n(\partial_nA_m,A^*_m)+2\int(A_m,A^*_m)$ with $E=x_n\partial_n$, along $P_b$ it is $\int(b_n\partial_nA_m,A^*_m)$; it sits between $\langle A_mA_n\rangle=2(-,-)^{-1}\delta_{mn}$ and $\langle A^*_m\,c\rangle=-(-,-)^{-1}ip_m$, each times $e^{-\varepsilon|p|^2}|p|^{-2}$ at $T=\infty$, and the insertion factors are $2$, $2i(l+q_c)_n$ and $ib_n(l+q_c)_n$. The special edge carries $e^{-\varepsilon|p|^2}-1$ times $(-,-)^{-1}$ and the inverse fibre pairing of $(A_m,A^*_m)$, $\tfrac14(A^{ij},A^*_{ij})$, $(\lambda,\lambda^*)$ and $(c,c^*)$: on $(\lambda_\alpha{}^i,\lambda^{*\beta}{}_j)$ the factor $-|l+q_\sigma|^2\delta_\alpha^\beta\delta^i_j(-,-)^{-1}$, on $(\bar\lambda_{\dot\alpha i},\bar\lambda^{*\dot\beta j})$ the opposite sign.
\item Bilinears. For $O=\tfrac1{16}\int\omega\,(w_1^{ij}A_{ij},w_2^{kl}A_{kl})$ the composite pair kernel along $\Delta$ is
\[
\Big[-\tfrac13\big((2E+2)A_{ij}\big)A_{kl}-\tfrac13A_{ij}\big((2E+2)A_{kl}\big)-\tfrac43A_{ij}A_{kl}\Big]\ \text{units},
\]
and along $P_b$ it is $-\tfrac13$ units times $b_m\partial_m(A_{ij}A_{kl})$, with the same contraction.
\item Two-site density. On adjacent pairs $(A_{ij},A_{kl})$, $(Z,\psi)$, $(\psi,Z)$, $(\psi,\psi)$ of a single trace the composite pair kernels give $-\tfrac43$ units (translation-invariant part along $\Delta$) and $-\tfrac16$ units times $b_m\partial_m$ (along $P_b$) on the same word, and $0$ on the exchanged word.
\item Free directions. The sector parts of $S^{(1)}(\Delta)$ are $\tfrac14\int\big((2E+2)A^{ij},A^*_{ij}\big)$, $-\int\big[\big((2E+3)\lambda_\alpha{}^i,\lambda^{*\alpha}{}_i\big)-\big((2E+3)\bar\lambda_{\dot\alpha i},\bar\lambda^{*\dot\alpha i}\big)\big]$ and $\int\big((2E+2)A_m,A^*_m\big)-\int\big(2E(c),c^*\big)$; those of $S^{(1)}(P_b)$ are $\tfrac14\int(b_m\partial_mA^{ij},A^*_{ij})$, $-\int\big[(b_m\partial_m\lambda_\alpha{}^i,\lambda^{*\alpha}{}_i)-(b_m\partial_m\bar\lambda_{\dot\alpha i},\bar\lambda^{*\dot\alpha i})\big]$ and $\int(b_m\partial_mA_n,A^*_n)-\int(b_m\partial_mc,c^*)$. The exact directions are $Q\int x_m(A_m,c^*)$ and $Q\int(\rho(x)\lambda^*,\lambda^*)=-2i\,Q\int x_m\big(\lambda^{*\alpha}{}_i,\sigma^m_{\alpha\dot\alpha}\bar\lambda^{*\dot\alpha i}\big)$. The further class of Proposition \ref{prop:c11-tc0} without the reflection is a multiple of $\int\big[(\lambda_\alpha{}^i,\lambda^{*\alpha}{}_i)+(\bar\lambda_{\dot\alpha i},\bar\lambda^{*\dot\alpha i})\big]$, the rotation $\lambda_\alpha{}^i\mapsto e^{i\varphi}\lambda_\alpha{}^i$, $\bar\lambda_{\dot\alpha i}\mapsto e^{-i\varphi}\bar\lambda_{\dot\alpha i}$ of the free gaugino.
\item Letter count. It counts $2$ for each letter $A_m$, $A_{ij}$, $\lambda_\alpha{}^i$, $\bar\lambda_{\dot\alpha i}$ and $1$ for each derivative $\partial_m$ on it; (b) reads: the composite pair kernels along $P_b$ sum to $-\tfrac13$ units times $b_m\partial_m$ on cohomology.
\item Families and spectra. $n_m=\sigma^m_{1\dot1}$, $D_n=n_m(\partial_m+[A_m,-])$, $\psi=\lambda_1{}^i$, and $F_{11}$ the self-dual field strength with both undotted indices $1$. In Srednicki's normalisation, $A_m=-igA^a_mT^a$, $A_I=-igA^a_IT^a$, $\lambda=-\tfrac{ig}{\sqrt2}\lambda^aT^a$, $\operatorname{tr}(T^aT^b)=\tfrac12\delta^{ab}$, the unit is $-\frac{g^2N}{32\pi^2}$ and $-\tfrac13$ units is $\frac{g^2N}{96\pi^2}$. On the families $D_1=-8$ units times the eigenvalues $E$ of $\sum_lH_{l,l+1}$, which is $\frac{g^2N}{4\pi^2}E=2\gamma$ with the one-loop anomalous dimension $\gamma=\frac{g^2N}{8\pi^2}E$: $\gamma=\frac{g^2N}{2\pi^2}h(S)$ on twist two, $\gamma=\frac{3g^2N}{4\pi^2}$ for $\operatorname{tr}(F_{11}F_{11})$ and the trace bilinear, the Konishi multiplet.
\item Odd generators. $Q_1(O)=\xi_\alpha{}^i\,Q_1(O)^\alpha{}_i+\bar\xi_{\dot\alpha i}\,\bar Q_1(O)^{\dot\alpha i}$ and $S_1(O)=\eta^\alpha{}_i\,S_1(O)_\alpha{}^i+\bar\eta^{\dot\alpha i}\,\bar S_1(O)_{\dot\alpha i}$ (Definition \ref{def:c8-components}); the Killing spinor of $\eta$ has the components $(\rho(x)\eta)_\alpha{}^i=-i\,x_m\sigma^m_{\alpha\dot\alpha}\bar\eta^{\dot\alpha i}$ and $(\rho(x)\eta)_{\dot\alpha i}=i\,x_m\eta^\alpha{}_i\sigma^m_{\alpha\dot\alpha}$ \eqref{eq:c8-pairing}, so on a family with $P_1\cong0$ the position part along $b$ of $S_1$ along $\eta$ is, on cohomology, $Q_1$ along $\xi_\alpha{}^i=-i\,b_m\sigma^m_{\alpha\dot\alpha}\bar\eta^{\dot\alpha i}$, $\bar\xi_{\dot\alpha i}=i\,b_m\eta^\alpha{}_i\sigma^m_{\alpha\dot\alpha}$. For $\xi_\alpha{}^i=u_\alpha z^i$, $(\rho(w)\xi)^\alpha{}_j=\varepsilon^{\alpha\beta}w_{ji}u_\beta z^i$, so $\rho(w)\xi=0$ exactly when $w_{ij}z^j=0$, and then $\iota_wA=\tfrac14w^{ij}A_{ij}$ is a combination of the components $A_{ij}z^j$: Proposition \ref{prop:c11-plane} is the vanishing of $Q_1$ on $\int\omega\,(m_1^iA_{ij}z^j,m_2^kA_{kl}z^l)$, $m_1,m_2\in W$.
\item Weinberg. On the Dirac column $(\lambda_\alpha{}^i,\bar\lambda^{\dot\alpha}{}_i)$, $\rho(v)=\gamma_W^\mu v_\mu$ gives Weinberg's form of each line, as in Remark \ref{rem:c10-twocomp}.
\end{itemize}
\end{remark}
